\subsection{Elliptic conventions and the exact Hodge splitting}

Throughout this subsection, $k\subset\Qb\subset\C$ is equipped with its indicated embedding. We use the degree-one motive of an elliptic curve, in the category of motives for absolute Hodge cycles generated by abelian varieties. The calculations of realizations below are also calculations of the ordinary algebraic de Rham and singular cohomology groups. Write
\[
 E/k:\quad y^2=4x^3-g_2x-g_3,\qquad
 \Delta=g_2^3-27g_3^2\ne0,
 \quad \omega=dx/y,\quad \xi=x\,dx/y.
\]
The classes $(\omega,\xi)$ form a $k$-basis of $H^1_{\dR}(E/k)$ with $F^1=k\omega$. The transcribed page 18:13 specifies a class in the additive quotient $\C/\Qb$, without choosing a unique second differential. Our scalar $h$ below is a representative relative to this explicit basis. This choice is essential when an exact numerical value, rather than algebraicity, is asserted.

Uniformize $E(\C)$ by $\C/\Lambda$, where
\[
 \Lambda=\Z\omega_1+\Z\omega_2,\quad
 \tau=\omega_2/\omega_1,\quad\operatorname{Im}\tau>0,
 \quad (x,y)=(\wp_\Lambda(z),\wp'_\Lambda(z)).
\]
The $\omega_i$ are full periods. Set $\eta_i=\zeta_\Lambda(z+\omega_i)-\zeta_\Lambda(z)$. Since $\zeta'=-\wp$, the periods of $\xi$ are $-\eta_i$. For the corresponding oriented homology basis, integration is the $\C$-linear map
\[
 I:H^1_{\dR}(E/k)\otimes_k\C
 \longrightarrow \operatorname{Hom}_{\Z}(H_1(E(\C),\Z),\C),
 \qquad I(v)(\gamma)=\int_\gamma v.
\]
With forms as rows and cycles as columns its matrix and determinant are
\begin{equation}\label{ell:periodmatrix}
 P=\begin{pmatrix}\omega_1&\omega_2\\-\eta_1&-\eta_2\end{pmatrix},
 \qquad \det P=\eta_1\omega_2-\eta_2\omega_1=2\pi i.
\end{equation}
This is Legendre's relation with full periods and the stated signs.

\begin{result}{A.4.1}{The scalar of the de Rham Hodge splitting}{Proved}
Writing $H^{0,1}=\C(\xi+h\omega)$, one has
\begin{align}
 h&=\frac{\eta_1\overline{\omega_2}-\eta_2\overline{\omega_1}}
 {\omega_1\overline{\omega_2}-\omega_2\overline{\omega_1}}\label{ell:hperiod}\\
 &=\frac{\eta_1}{\omega_1}-\frac{\pi}{\omega_1^2\operatorname{Im}\tau}
 =\frac{\pi^2}{3\omega_1^2}E_2^*(\tau).\label{ell:he2}
\end{align}
For an algebraic adapted change of basis $\omega'=r\omega$, $\xi'=s\xi+t\omega$, with $r,s\in\Qb^\times$, $t\in\Qb$, the parameter is $h'=(sh-t)/r$. In particular, the property $h\in\Qb$ is independent of the adapted algebraic basis.
\end{result}
\begin{proof}
The conjugate of the holomorphic class has periods $\overline{\omega_i}$. Write it as $a\omega+b\xi$; thus $a\omega_i-b\eta_i=\overline{\omega_i}$. Put $A=\operatorname{Im}(\overline{\omega_1}\omega_2)=|\omega_1|^2\operatorname{Im}\tau>0$. Solving the two equations by \eqref{ell:periodmatrix} gives $b=-A/\pi\ne0$ and $a/b$ equal to \eqref{ell:hperiod}. Substituting $\omega_2=\tau\omega_1$ and $\eta_2=\tau\eta_1-2\pi i/\omega_1$ gives the first expression in \eqref{ell:he2}. The identity $\eta_1\omega_1=\pi^2 E_2(\tau)/3$ gives the second. Finally,
$\xi'+h'\omega'=s(\xi+h\omega)$ if and only if $t+rh'=sh$, which proves the asserted basis transformation.
\end{proof}

Here the Fourier and Weierstrass normalizations are
\begin{align*}
 q&=e^{2\pi i\tau},& E_2&=1-24\sum_{n\ge1}\sigma_1(n)q^n,
 &E_2^*&=E_2-\frac3{\pi\operatorname{Im}\tau},\\
 E_4&=1+240\sum_{n\ge1}\sigma_3(n)q^n,
 &E_6&=1-504\sum_{n\ge1}\sigma_5(n)q^n,\\
 g_2&=\frac{4\pi^4}{3\omega_1^4}E_4,
 &g_3&=\frac{8\pi^6}{27\omega_1^6}E_6.
\end{align*}
Katz's $S(E,\omega)$ equals $-12h$. His Fourier choice $\omega_1=2\pi i$ consequently gives $S=E_2^*$. This is the normalization of \cite[Lemma 1.3.6 and (1.3.7), p.467; (1.3.10), p.468]{EllKatz1976}, with the algebraic de Rham basis in his \S1.2. Replacing full periods by half-periods requires changing Legendre's relation as well.

\subsection{The filtration field and the bigrading field}

\begin{result}{A.3.1}{The two elliptic flag fields}{Proved, using Schneider's theorem}
In the rational Betti basis dual to the cycles used above,
\[
 F^1=\C(1,\tau),\qquad H^{0,1}=\C(1,\overline\tau),
 \qquad \Q(\varphi)=\Q(\tau),\quad
 \Q(\psi)=\Q(\tau,\overline\tau).
\]
For $E/\Qb$, the filtration field has transcendence degree zero in the CM case and one in the non-CM case. The bigrading field has transcendence degree zero in the CM case; in the non-CM case its transcendence degree lies between one and two.
\end{result}
\begin{proof}
Integration of $\omega$ gives $(\omega_1,\omega_2)$, and integration of its complex conjugate gives $(\overline{\omega_1},\overline{\omega_2})$. Both first coordinates are nonzero. The affine coordinate of each line in $\mathbf P^1$ is therefore exactly its period ratio. The ordered-pair chart for the two transverse lines has coordinates $\tau,\overline\tau$, proving the field equalities. Schneider's original theorem, \cite[\S1, result IIc, p.5]{EllSchneider1937}, says that an algebraic $j(\tau)$ permits algebraic $\tau$ only when $\tau$ is imaginary quadratic. Conversely an imaginary-quadratic period ratio gives complex multiplication: multiplication by a suitable nonrational integral element of $\Q(\tau)$ preserves a finite-index sublattice and, after multiplying the element by an integer, preserves $\Z+\Z\tau$. Thus the two cases have the stated transcendence degrees. In the CM case $\overline\tau$ is also algebraic; the remaining bounds follow from the displayed field equalities.
\end{proof}

\begin{result}{A.3.2}{Curves isogenous to their complex conjugates}{Proved}
If $E/\Qb$ is isogenous to $\overline E$, then $\overline\tau\in\Q(\tau)$, so $\Q(\psi)=\Q(\varphi)$. In the non-CM case their common transcendence degree over $\Q$ is one. Moreover
\[
 h\in\Qb(\omega_1,\omega_2,\eta_1,\eta_2).
\]
This applies to every curve with a real model and also to curves isogenous to their conjugates without having a real model.
\end{result}
\begin{proof}
A holomorphic isogeny is induced on universal covers by a nonzero scalar $\lambda$ with $\lambda\Lambda\subset\overline\Lambda$. Therefore there are integers $a,b,c,d$ with $ad-bc\ne0$ such that
\[
 \lambda\omega_1=a\overline{\omega_1}+b\overline{\omega_2},
 \qquad \lambda\omega_2=c\overline{\omega_1}+d\overline{\omega_2}.
\]
It follows that
\[
 \tau=\frac{c+d\overline\tau}{a+b\overline\tau},
 \qquad \overline\tau=\frac{c-a\tau}{b\tau-d}.
\]
The last denominator cannot vanish: its vanishing would also force $c-a\tau=0$, contradicting the nonzero determinant. Conversely, a rational fractional-linear relation of this form with nonzero determinant gives commensurable lattices after scaling, hence an analytic isogeny. The isogeny descends to $\Qb$: its finite kernel consists of torsion points of the algebraic curve, its quotient is defined over $\Qb$, and a curve with the target $j$-invariant is isomorphic to the target over $\Qb$.

For an isogeny over $\Qb$, pullback of an invariant differential gives $f^*\overline\omega=\lambda\omega$ with $\lambda\in\Qb^\times$. Inverting the displayed integral matrix over $\Q$ puts $\overline{\omega_1},\overline{\omega_2}$ in $\Qb(\omega_1,\omega_2)$. Formula \eqref{ell:hperiod} proves the last inclusion. A real model gives an isomorphism to the conjugate, hence is a special case.
\end{proof}

\begin{result}{A.3.3}{An explicit nonreal example with bigrading transcendence degree one}{Proved}
The curve
\[
 E_7/\Q(i):\quad Y^2=X^3+2X^2+(\tfrac12+i)X
\]
is non-CM, has no real model, and is isogenous to its conjugate. Its bigrading field has transcendence degree one.
\end{result}
\begin{proof}
For $a=2$, $b=1/2+i$, the degree-two isogeny
\[
 (X,Y)\longmapsto\bigl(X+a+b/X,\,Y(1-b/X^2)\bigr)
\]
has target $Y'^2=X'^3-4X'^2+(2-4i)X'$. The formula follows by substitution into the target equation; the map extends at its poles to the smooth projective curves, preserves the identity, and has degree two. If $u^2=-2$, scaling $\overline E_7$ by $(X,Y)\mapsto(u^2X,u^3Y)$ gives exactly this target because $u^2\overline a=-4$ and $u^4\overline b=2-4i$.

Direct evaluation gives
\[
 j(E_7)=\frac{-56512+38144i}{25},\qquad
 N_{\Q(i)/\Q}(j(E_7))=\frac{929714176}{125}.
\]
The nonintegral norm proves that $j(E_7)$ is not an algebraic integer, whereas a CM $j$-invariant is integral; an opened proof of the latter is \cite[Theorem 12.14 and Corollaries 12.15--12.16, pp.119--120]{EllMilneMF}. Its nonreal $j$ precludes a real model because the invariant of a real equation is real and is invariant under complex isomorphism. Apply A.3.2. The short-Weierstrass invariants are $g_2=10/3-4i$, $g_3=-28/27+8i/3$. The symbolic program \texttt{check\_conjugate\_examples.py} verifies all these identities.
\end{proof}

\begin{result}{A.3.4}{An explicit curve not isogenous to its conjugate}{Proved}
Let $j_0=1/(2+i)$, $c=j_0/(j_0-1728)$, and
\[
 E_8/\Q(i):\quad y^2=4x^3-3cx-c.
\]
Then $j(E_8)=j_0$. This curve is non-CM, has no real model, and is not isogenous to $\overline E_8$.
\end{result}
\begin{proof}
Substitution in $j=1728g_2^3/(g_2^3-27g_3^2)$ gives $j_0$; the discriminant is nonzero because $c\ne0,1$. At the Gaussian prime $v=(2+i)$ one has $v(j_0)=-1$, but $v(\overline{j_0})=0$: in the residue field $\mathbf F_5$, $i=3$ and $2-i=4\ne0$. Nonintegrality gives non-CM, and nonreality gives the absence of a real model as in A.3.3.

Any isogeny between complex elliptic curves can be replaced by a cyclic isogeny between the same isomorphism classes. Indeed a kernel of type $\Z/m\oplus\Z/n$, $m\mid n$, contains $E[m]$: its $m$-torsion has order $m^2$ and equals $E[m]$. Factoring through $[m]$ leaves a cyclic kernel of order $n/m$. A cyclic $N$-isogeny between $E_8$ and its conjugate would give $\Phi_N(j_0,\overline{j_0})=0$. The modular polynomial has integer coefficients and is monic in either variable, by \cite[Theorem 6.1 and its proof, pp.82--83]{EllMilneMF}. Consequently $j_0$ would be integral over the valuation ring containing $\overline{j_0}$, contradicting $v(j_0)=-1$. For $N=1$ the same conclusion follows from equality of the $j$-invariants. The defining identities are checked symbolically in the cited program.
\end{proof}

\paragraph{Unresolved gate A.3.}
For $E$ and $\overline E$ nonisogenous, the remaining assertion is algebraic independence of $\tau,\overline\tau$. A.3.2 excludes rational fractional-linear relations in this case; its proof does not exclude polynomial relations of higher degree. For the explicit $E_8$, the bounds proved here remain
$1\le\trdeg_{\Q}\Q(\tau,\overline\tau)\le2$. No opened theorem in this run settles this gate.

\subsection{CM algebraicity, its exact value, and the CM period field}

\begin{result}{A.4.2}{Algebraic Hodge splitting for CM curves}{Known, with derivation}
For a CM elliptic curve $E/k$ in the stated Weierstrass normalization, $h\in k$. In particular its Hodge splitting is defined over $k$.
\end{result}
\begin{proof}
First choose a nonreal CM endomorphism $u$, defined over a finite extension of $k$. On the holomorphic and antiholomorphic lines its eigenvalues are $u$ and $\overline u$, respectively. In the de Rham basis its action is
\[
 u^*\omega=u\omega,\qquad u^*\xi=a\omega+\overline u\,\xi,
 \qquad a\in\Qb.
\]
Applying it to $\xi+h\omega$ gives $a+uh=\overline u h$, hence $h=a/(\overline u-u)\in\Qb$. This is the eigenspace mechanism of \cite[Lemma 4.0.7, p.501]{EllKatz1976}. To descend to the coefficient field itself, Masser's Theorem A1 gives $\Psi(\tau)\in\Q(j(\tau))$ for his function
$\Psi=3E_4E_2^*/(2E_6)$ at an imaginary-quadratic point outside the orbit of $i$; see \cite[Appendix One, Theorem A1 and (100), p.114]{EllMasser1975}. If $j\ne0,1728$, the explicit normalizations give $h=(g_3/g_2)\Psi(\tau)\in k$. At $j=0,1728$, the modular stabilizers force $E_2^*=0$ and therefore $h=0\in k$. This also handles the two exceptional denominators in the displayed formula.
\end{proof}

\begin{result}{A.4.3}{The full CM period field}{Known, with exact field calculation}
For a CM elliptic curve over $\Qb$, let
\[
 L=\Qb(\omega_1,\omega_2,\eta_1,\eta_2).
\]
Then
\[
 L=\Qb(\omega_1,\pi/\omega_1)=\Qb(\omega_1,\pi),
 \qquad \trdeg_{\Qb}L=2.
\]
For the degree-one motive, the CM motivic torus has dimension two; this proves the period-conjecture dimension in this case. The same transcendence degree holds after adjoining the periods of the conjugate curve.
\end{result}
\begin{proof}
The numbers $\tau$ and $\operatorname{Im}\tau$ are algebraic and nonzero in the CM case, and A.4.2 gives algebraic $h$. Formula \eqref{ell:he2} gives
\[
 \omega_2=\tau\omega_1,\qquad
 \eta_1=h\omega_1+\frac{\pi}{\omega_1\operatorname{Im}\tau},
 \qquad \eta_2=\tau\eta_1-\frac{2\pi i}{\omega_1}.
\]
Thus $L\subset\Qb(\omega_1,\pi/\omega_1)$. Legendre's relation places $\pi$ in $L$, proving the reverse inclusion. Chudnovsky's theorem gives transcendence degree at least two for the four periods and quasiperiods of an elliptic curve with algebraic invariants; we use the original author's \cite[\S2A, Theorem 1, p.343, with the algebraic-invariant convention on p.339]{EllChudnovsky1978}. The displayed field has at most two generators, so equality follows. This also proves algebraic independence of $\pi$ and $\omega_1$.

The CM action confines the Mumford--Tate group to the two-dimensional torus $\operatorname{Res}_{K/\Q}\Gm$, where $K$ is the CM field. Over $\C$ the Hodge cocharacter and its Galois conjugate are the two coordinate cocharacters, so they generate that torus. In the category of abelian motives for absolute Hodge cycles this is the motivic group, by the absolute-Hodge theorem used for the category in the paper.

For completeness, the conjugate CM lattice is commensurable after scaling with the original: both normalized lattices span the same imaginary-quadratic field over $\Q$. The resulting isogeny descends to $\Qb$ by the kernel argument in A.3.2. Its matrix on de Rham cohomology is invertible over $\Qb$, and its matrix on rational homology is invertible over $\Q$. Integration therefore expresses all periods and quasiperiods of the conjugate curve as $\Qb$-linear combinations of those of $E$. The compositum has the same field $L$. For any algebraic base field $k$ over which the data are defined, algebraic extension to $\Qb$ does not change these transcendence degrees.
\end{proof}

\begin{result}{A.4.4}{An exact nonzero CM splitting}{Proved}
For $E:y^2=4x^3-30x-28$ one has $j=8000$, CM by $\Q(\sqrt{-2})$, and $h=1/2$ in the basis $(dx/y,x\,dx/y)$.
\end{result}
\begin{proof}
Put $X=x+2$, $Y=y/2$. The equation becomes $E_0:Y^2=X^3-6X^2+9X/2$. The formulas
\[
 X'=X-6+\frac9{2X},\qquad Y'=Y\left(1-\frac9{2X^2}\right)
\]
give a normalized degree-two isogeny $\varphi:E_0\to E_1$, where $E_1:Y'^2=X'^3+12X'^2+18X'$. Substitution verifies the target equation, and pullback sends $dX'/(2Y')$ to $dX/(2Y)$. The induced function-field extension has degree two, as is also seen from the quadratic equation for $X$ in terms of $X'$. The rational map extends to the smooth projective curves and carries the origin to the origin.

For $u=\sqrt{-2}$, the map $s:E_0\to E_1$, $(X,Y)\mapsto(u^2X,u^3Y)$, is an isomorphism. Hence $\alpha=s^{-1}\varphi$ is a degree-two endomorphism with $\alpha^*\omega=u\omega$. Its nonrational eigenvalue proves CM. Moreover
\[
 d(Y/X)=\left(X-\frac9{2X}\right)\omega,
 \qquad \alpha^*\xi/u=(-X'/2-2)\omega
 =(-X+1)\omega+\tfrac12d(Y/X).
\]
Because $\xi=(X-2)\omega$, this gives $\alpha^*\xi=-u\xi-u\omega$ in de Rham cohomology. The eigenspace calculation in A.4.2 yields $h=(-u)/(-2u)=1/2$. The discriminant is $5832$, giving $j=1728\cdot30^3/5832=8000$. The program \texttt{verify\_cm\_isogeny.py} checks every rational-function identity exactly.
\end{proof}

\begin{result}{A.4.5}{Masser's explicit CM formula in these conventions}{Known, with normalization checked}
For imaginary-quadratic $\tau$ with $j=j(\tau)\ne0,1728$, define $\Psi=3E_4E_2^*/(2E_6)$. Then
\[
 h=(g_3/g_2)\Psi,\qquad
 \Psi=9j\gamma+\frac{3(7j-6912)}{2(j-1728)},
\]
where $\gamma\in\Q(j)$ is the modular-polynomial expression specified below. At $j=0,1728$, one has $h=0$.
\end{result}
\begin{proof}[Derivation and exact locator]
Choose a primitive equation $C\tau^2+B\tau+A=0$, $C>0$, and set
\[
 \mu=\begin{cases}
 \begin{pmatrix}-B&-2A\\2C&B\end{pmatrix},&B\text{ odd},\\[2mm]
 \begin{pmatrix}-B/2&-A\\C&B/2\end{pmatrix},&B\text{ even},
 \end{cases}\qquad D=\det\mu.
\]
Expand $\Phi_D(j+X,j+Y)=\sum\beta_{mn}X^mY^n$. In the ordinary case,
\[
 \beta=\beta_{10}=\beta_{01}\ne0,\qquad
 \gamma=(\beta_{20}-\beta_{11}+\beta_{02})/\beta.
\]
In Masser's exceptional case, $B$ is odd and $D=3d^2$ for an integer $d>1$, and
\[
 \beta=\beta_{30}=\beta_{03}\ne0,\quad \beta_{21}=\beta_{12}=0,
 \quad \gamma=(\beta_{40}-\beta_{31}+\beta_{22}-\beta_{13}+\beta_{04})/\beta.
\]
These are \cite[Lemmas A1--A2, pp.116--118, and (105)--(106), pp.118--119]{EllMasser1975}. The equality for $h$ follows by dividing the exact expressions for $g_3,g_2$ above. The application of (106) is restricted here to $j\ne0,1728$: its displayed derivation divides by $j'(\tau)$. At the elliptic points one instead uses $E_2^*=0$. Naive substitution of $j=0$ in the rational expression would give a spurious value. The table on Masser's p.121 gives $\Psi(i\sqrt2)=15/28$, and $(28/30)(15/28)=1/2$, independently checking the normalization in A.4.4.
\end{proof}

\paragraph{Two transcription hazards in the later literature.}
For numerical formulas, use Masser's original (100). The abbreviated description of his function in \cite[p.774]{EllSpence2018} omits its factor $3/2$; the field-inclusion argument is unaffected. On that same printed page, the exceptional expression displays the denominator $\beta_{01}$, whereas Masser's exceptional denominator is $\beta_{03}$ and $\beta_{01}$ vanishes. Both printed pages were visually checked; this is not an OCR substitution. The statements used here are taken from Masser's original formula.

\subsection{The zero value and isogeny transport}

\begin{result}{A.4.6}{Complete classification of the zero scalar}{Proved}
For any complex elliptic curve in the short-Weierstrass normalization fixed above,
\[
 h=0\quad\Longleftrightarrow\quad j(E)\in\{0,1728\}.
\]
\end{result}
\begin{proof}
By \eqref{ell:he2}, it suffices to locate the zeros of $E_2^*$. Its weight-two transformation law reduces the question to $|x|\le1/2$, $|\tau|\ge1$, for $\tau=x+iy$. Thus $y\ge\sqrt3/2$ and $r=e^{-2\pi y}<1/200$. For $0<x<1/2$, use $\sigma_1(n)\le n^2$ and $|\sin(n\theta)|\le n\sin\theta$ for $0<\theta<\pi$ to obtain
\[
 \left|\sum_{n\ge2}\sigma_1(n)r^n\sin(2\pi nx)\right|
 \le\sin(2\pi x)\sum_{n\ge2}n^3r^n
 <r\sin(2\pi x).
\]
Indeed
\[
 S_3(r)=\sum_{n\ge1}n^3r^n=\frac{r(1+4r+r^2)}{(1-r)^4},
 \qquad S_3(1/200)/(1/200)<2.
\]
The first Fourier term therefore gives a nonzero imaginary part. Reflection deals with $-1/2<x<0$. A zero must lie on $x=0$ or $x=\pm1/2$.

For $x=0$, $y\ge1$ and
\[
 \frac d{dy}E_2^*(iy)=48\pi\sum_{n\ge1}n\sigma_1(n)e^{-2\pi ny}
 +\frac3{\pi y^2}>0.
\]
For $x=1/2$, $y\ge y_0=\sqrt3/2$, and $y^2e^{-2\pi ny}$ decreases on this ray. Consequently
\[
 y^2\frac d{dy}E_2^*(1/2+iy)
 \ge\frac3\pi-36\pi S_3(e^{-\pi\sqrt3})
 >\frac3\pi-\frac\pi5>0.
\]
Here $S_3(1/200)=8160200/1568239201<1/180$ and $\pi^2<15$. Modular stabilizers force zero at $i$ and at $\rho+1$, where $\rho=e^{2\pi i/3}$. Strict monotonicity proves uniqueness on each ray; reflection and translation finish the fundamental domain. The orbits have $j=1728$ and $j=0$.

The analytic lemma and proof are also in the opened primary source \cite[Lemma 8.1, printed pp.25--26]{EllSebbarBesrour2026}. Every estimate just used was independently checked with rational bounds by \texttt{check\_zero\_bounds.py}. This assertion concerns the scalar in the fixed short-Weierstrass basis. Its class in $\C/\Qb$ vanishes for all CM curves, including A.4.4 where the scalar is nonzero.
\end{proof}

\begin{result}{A.4.7}{Algebraic isogeny transport of the splitting scalar}{Proved}
Let $f:E\to E'$ be a degree-$n$ isogeny over $\Qb$ between short-Weierstrass models, and write
\[
 f^*\omega'=\lambda\omega,\qquad f^*\xi'=a\omega+b\xi.
\]
Then $\lambda,a,b\in\Qb$, $\lambda b=n$, and
\[
 h_{E'}=\frac{n}{\lambda^2}h_E-\frac a\lambda.
\]
Hence algebraicity of the splitting is constant on an algebraic isogeny class.
\end{result}
\begin{proof}
Pullback on algebraic de Rham cohomology has algebraic coefficients and preserves $F^1$. The cup-product pairing is normalized by $(\omega,\xi)_{\dR}=1$, as in Katz \S1.2. Pullback multiplies this pairing by $n$, so its determinant is $\lambda b=n$ and both factors are nonzero. Preservation of $H^{0,1}$ gives
\[
 a\omega+b\xi+\lambda h_{E'}\omega=b(\xi+h_E\omega).
\]
Solving for $h_{E'}$ proves the formula and the assertion.
\end{proof}

\subsection{The almost-holomorphic formulation and the exact remaining gate}

\begin{result}{A.4.8}{Translation into the function $\chi^*$}{Proved}
For $j\ne0,1728$, define
\[
 \chi^*(\tau)=\frac{1728E_2^*(\tau)E_4(\tau)E_6(\tau)}
 {E_4(\tau)^3-E_6(\tau)^2}.
\]
Then
\[
 \chi^*=\tfrac23(j-1728)\Psi,\qquad
 h=\frac{3g_3}{2g_2(j-1728)}\chi^*.
\]
For algebraic $g_2,g_3\ne0$, algebraicity of $h$ is equivalent to algebraicity of $\chi^*(\tau)$.
\end{result}
\begin{proof}
Substitute $j-1728=1728E_6^2/(E_4^3-E_6^2)$ and the definition of $\Psi$. The coefficient in the second equality is a nonzero algebraic number under the stated hypotheses, proving the equivalence.
\end{proof}

Spence's \cite[Proposition 5.1, p.773]{EllSpence2018} proves $\Q(\chi^*(\tau))\subseteq\Q(j(\tau))$ at quadratic $\tau$. His Theorem 5.8, p.776, proves finiteness of maximal $(j,\chi^*)$-special subvarieties. Definition 3.6, pp.755--756, defines their constant factors using quadratic $\tau$. Applied to an arbitrary algebraic singleton, the finiteness theorem allows either a special singleton or no special subvariety at all. This specialization therefore does not prove that every algebraic pair $(j(\tau),\chi^*(\tau))$ is CM. The calculation states exactly what that application yields; it does not exclude a further argument extending the methods.

\paragraph{Unresolved gate A.4(c).}
The converse in Katz's \cite[Question 4.0.8, p.501]{EllKatz1976} is not established by the results above. Its concrete real test case here is
\[
 E_* /\Q:\quad y^2=4x^3-4x-1,
 \qquad \Delta=37,\quad j=110592/37.
\]
Nonintegrality of $j$ proves non-CM. The cubic has three distinct real roots, so choose $\omega_1>0$ real and $\omega_2\in i\R_{>0}$. Reality of the Weierstrass series gives $\eta_1\in\R$, $\eta_2\in i\R$. Formula \eqref{ell:hperiod} becomes
\[
 h_* =\frac{\eta_1\omega_2+\eta_2\omega_1}{2\omega_1\omega_2}.
\]
For any $a\in\Qb$, there are exact equivalences
\begin{equation}\label{ell:exactgate}
 \begin{split}
 h_*=a
 &\quad\Longleftrightarrow\quad
 \eta_1\omega_2+\eta_2\omega_1-2a\omega_1\omega_2=0\\
 &\quad\Longleftrightarrow\quad (\eta_1-a\omega_1)\omega_2=\pi i.
 \end{split}
\end{equation}
The second equivalence follows from Legendre's relation with the signs in \eqref{ell:periodmatrix}. The unresolved assertion is that \eqref{ell:exactgate} fails for every algebraic $a$. A.4.6 excludes $a=0$; it does not exclude the other algebraic values. The computed approximation is
\[
 h_*=0.113219334661746861539065791000647532558647974361219962657816837416979\ldots.
\]
No algebraicity inference is drawn from it. For curves without a real model, formula \eqref{ell:hperiod} is the exact obstruction with the conjugate periods retained; A.3.2 proves their elimination when the curve is isogenous to its conjugate.

\subsection{Tests of the classical transcendence theorems}

Chudnovsky's \cite[\S2A, Theorems 1 and 3, p.343]{EllChudnovsky1978} gives, respectively, transcendence degree at least two for the four periods and quasiperiods, and algebraic independence of $\pi/\omega$ and $\eta/\omega$. Here is an explicit algebraic test of the proposed contradiction from the first assertion and linear period independence. Write
\[
 x=\omega_1,\quad y=\omega_2,\quad
 u=\eta_1-a\omega_1,\quad v=\eta_2-a\omega_2.
\]
The putative relation is $uy+vx=0$. On the open set $xyu\ne0$, the algebraic variety $v=-uy/x$ has function field $\Qb(x,y,u)$ of transcendence degree three, and its four coordinate functions are $\Qb$-linearly independent. Indeed multiplying a proposed constant linear relation by $x$ gives a linear combination of the distinct monomials $x^2,xy,xu,uy$. Its determinant $uy-vx=2uy$ is nonzero. This is an algebraic compatibility witness for the displayed relation and those lower bounds, not a tuple of genuine periods and not a counterexample to a stronger period theorem.

Nesterenko's \cite[Theorem 1, p.1320]{EllNesterenko1996} gives transcendence degree at least three for $q,E_2(q),E_4(q),E_6(q)$ when $0<|q|<1$. His Corollary 3, p.1321, gives algebraic independence of
\[
 q,\qquad \omega_1/\pi,\qquad \eta_1/\pi
\]
for algebraic invariants. With $t=\omega_1/\pi$, the equality $h=a\in\Qb$ in \eqref{ell:he2} gives
\[
 \eta_1/\pi=at+\frac1{t\pi\operatorname{Im}\tau},
 \quad
 \Qb(q,t,\eta_1/\pi)=\Qb(q,t,\pi\operatorname{Im}\tau).
\]
Nesterenko's theorem then gives transcendence degree three for the generators on the right. It does not contradict this field equality. This checks the specialization of the exact theorem, with its normalizations retained.

The original Schneider source was obtained as the G\"ottingen scan and read at \S1, result IIc, p.5. The complete original Chudnovsky congress contribution was obtained from the official IMU volume and read at pp.339 and 343. The 2026 theorem of \cite[Theorem 1]{EllSenthilKumar2026} requires $u_1,u_2,\omega$ to be $\Q$-linearly independent and $(\Z u_1+\Z u_2)\cap\Lambda=0$; its at-least-two lower bound cannot be specialized by taking the two extra logarithms in the period lattice. Waldschmidt's \cite[\S3.3, p.31]{EllWaldschmidt2004} explicitly records Katz's question and points to Bertrand's nonholomorphic conjecture. This supplies historical status evidence, not proof of its present status. The referenced Bertrand chapter's metadata was opened, but its full text was not obtained and no unseen theorem from it is invoked.

\subsection{Reproducibility of the elliptic computations}

The exact-input programs accompany this section. The isogeny and conjugate examples use symbolic rational-function arithmetic. The zero estimates use rational inequalities, including $S_3(1/200)<1/180$ and a rational Taylor lower bound proving $e^{\pi\sqrt3}>200$. The numerical program evaluates $(g_2,g_3)=(4,0),(30,28),(4,1)$ at 80 and 160 decimal digits.

For real roots $e_1>e_2>e_3$, put $d=e_1-e_3$, $m=(e_2-e_3)/d$. It independently computes
\begin{align*}
 \omega_1&=2K(m)/\sqrt d,&
 \omega_2&=2iK(1-m)/\sqrt d,\\
 \eta_1&=2\sqrt d\,\mathrm E(m)-e_1\omega_1,&
 \eta_2&=-2i\sqrt d\,\mathrm E(1-m)-e_3\omega_2.
\end{align*}
Here $K$ and $\mathrm E$ are the complete elliptic integrals. In particular $\eta_2$ is not defined by forcing Legendre's relation. The program compares these values with the Fourier formula for $\eta_1$, verifies Legendre's relation, compares both formulas for $h$, and reconstructs $g_2,g_3$. Residuals in the 160-digit run are of order $10^{-160}$.

For a Fourier sum with divisor weight $s$, coefficient $c$, $N$ retained terms, $r=|q|$, and $R=r((N+2)/(N+1))^{s+1}<1$, the omitted tail is bounded by
\[
 |c|\frac{(N+1)^{s+1}r^{N+1}}{1-R},
\]
using $\sigma_s(n)\le n^{s+1}$. This bound is rigorous for exact $r$. The implementation evaluates the integrals and $r$ in high-precision floating arithmetic, so the output is a consistency check rather than a full interval certificate for rounding error. The programs, inputs, precision, and saved outputs are all included; no transcendence conclusion rests on numerical recognition.
\subsubsection{Exact modern reformulations and a corrected theorem}

\begin{result}{A.4.10}{A coefficient that detects the Hodge splitting}{Proved}
Use the period normalization of the optimal modular elliptic curve $E/\Q$ in \cite{EllAGOR2015}. Write its normalized newform and Eichler integral as
\[
 F_E(q)=\sum_{n\ge1}a_nq^n,\quad a_1=1,
 \qquad \mathcal E_E(q)=\sum_{n\ge1}\frac{a_n}{n}q^n.
\]
For the short-Weierstrass basis $\omega=dz$, $\xi=\wp(z)dz$, their regularized lattice sum $S(\Lambda)$ equals $h$. Their holomorphic Weierstrass mock modular form satisfies
\[
 \widehat{\mathfrak Z}_E^+(q)
 =\zeta_\Lambda(\mathcal E_E(q))-h\mathcal E_E(q),
 \qquad
 [q]\widehat{\mathfrak Z}_E^+
 =\frac{a_2^2}{4}-\frac{a_3}{3}-h.
\]
Thus $h\in\Qb$ is equivalent to algebraicity of this one coefficient, and to algebraicity of all the coefficients of $\widehat{\mathfrak Z}_E^+$.
\end{result}
\begin{proof}
Equation (2.2) of \cite{EllAGOR2015}, printed preprint p.9, states that
$\zeta_\Lambda(z)-S(\Lambda)z-(\pi/A)\bar z$ is $\Lambda$-periodic, where $A=\operatorname{Im}(\overline{\omega_1}\omega_2)$. Its period increments are therefore
$\eta_i-S(\Lambda)\omega_i-(\pi/A)\overline{\omega_i}=0$.
The linear system used in A.4.1 gives the same equations with $h$ in place of $S(\Lambda)$, proving their equality. In this algebraic differential normalization $g_2,g_3\in\Q$; all Laurent coefficients of $\zeta_\Lambda$ are rational, by the differential equation for $\wp$ and $\zeta'=-\wp$. The expansion
\[
 \mathcal E_E(q)=q+\frac{a_2}{2}q^2+\frac{a_3}{3}q^3+O(q^4)
\]
gives $[q](1/\mathcal E_E)=a_2^2/4-a_3/3$. Since $\zeta(z)=z^{-1}+O(z^3)$, this proves the formula. Every other coefficient is a rational number minus a rational multiple of $h$, proving the equivalences.
\end{proof}

For clarity, this is an exact reformulation of the unresolved arithmetic assertion. Remark~6 of \cite{EllAGOR2015}, printed p.5, records non-CM transcendence of $S(\Lambda)$ as an expectation. For $37\mathrm a1$, given by $y^2+y=x^3-x$, one has $Y=2y+1$, hence $Y^2=4x^3-4x+1$, and $(a_2,a_3)=(-2,-3)$. The displayed coefficient is $2-h$. Its value $2.113219\ldots$ agrees with A.4.1 and the algebraic twist by $i$ of the model in the explicit unresolved gate.

There is also a direct comparison with modern single-valued periods. For a real algebraic model and adapted basis, let $P_c=P^{\mathsf T}$ be the integration matrix with cycles as rows. The first column of
\[
 S=P_c^{-1}\overline{P_c}=(s_{ij})
\]
represents the conjugate holomorphic class, so $s_{21}=-A/\pi\ne0$ and $h=s_{11}/s_{21}$. Fonseca's Theorem~0.4 and Proposition~8.1 \cite{EllFonseca2020}, printed pp.6 and28--29, express negative-index Poincar\'e coefficients, in a rational symplectic modular basis $([f],[g])$ with $\dim S_{k+2}(\Gamma_0(N))=1$, as
\[
 a_n(P_{-m,k+2,N})
 =\frac{k!}{m^{k+1}}a_m(f)a_n(f)\frac{s_{11}}{s_{21}}+r_{m,n},
 \qquad r_{m,n}\in\Q.
\]
In weight two, $m=n=1$ detects the same algebraicity property. His $X_0(11)$ example uses $\xi_{\rm gen}=(\wp+1/3)dz$, so $h_{\rm short}=h_{\rm gen}+1/3$ and $a_1(P_{-1,2,11})=h_{\rm gen}-1$; the basis-change formula fixes this translation exactly.

\paragraph{Corrigendum audit.}
The original Bruinier--Ono--Rhoades Theorem~1.4, journal p.677, was replaced in their 2009 corrigendum \cite{EllBORCorr2009}. Its valid conclusion, for a prime $p\nmid N$ with vanishing Hecke eigenvalue $a_p$, gives algebraic holomorphic coefficients only at indices $n$ with $\operatorname{ord}_p(n)$ odd. In weight two, the recurrence
$a_{p^r}=a_p a_{p^{r-1}}-p a_{p^{r-2}}$
then gives $a_{p^{2r+1}}=0$, and multiplicativity gives $a_n=0$ at all of these indices. Accordingly, the coefficient of $h$ in the preceding formula vanishes there. The detecting index $n=1$ is not covered. The superseded unrestricted version cannot be used to claim either a counterexample or a solution of the gate.

For a further convention check, Brown--Fonseca \cite{EllBrownFonseca2025}, Lemma~6.3, equations (6.8)--(6.9), printed p.34, use $S=\overline{P_c}^{-1}P_c$ and $\rho=S_{11}/S_{21}$; their conjugate holomorphic class is $\xi+\bar\rho\omega$. Thus $h=\bar\rho$ in that convention. Lemma~6.4 proves the CM direction. Proposition~7.3 expresses its associated meromorphic modular forms in powers of $\rho$ and a cusp-form term; Proposition~7.5 specializes at CM points. The Gross--Zagier extension in Theorem~8.14 explicitly assumes both marked points CM. These formulas give exact comparisons but do not establish the non-CM converse.

\subsubsection{Exact tensor and functional-transcendence tests}

\begin{result}{A.4.11}{The projector and the rank-three tensor of the splitting}{Proved}
Let $e$ be projection onto $H^{0,1}$ along $H^{1,0}$. Its de Rham and Betti matrices are
\[
 e_{\dR}=\begin{pmatrix}0&h\\0&1\end{pmatrix},\qquad
 e_B=\frac1{\bar\tau-\tau}
 \begin{pmatrix}-\tau&1\\-\tau\bar\tau&\bar\tau\end{pmatrix}.
\]
In particular,
\[
 \Qb(\text{entries of }e_B)=\Qb(\tau,\bar\tau),
\]
and $e_B$ has algebraic entries exactly when $E$ has CM.
For a real model with real Frobenius $F=\operatorname{diag}(1,-1)$, put
$P_c=\left(\begin{smallmatrix}x&p\\y&q\end{smallmatrix}\right)$ and $D=\det P_c$. The inverse comparison image of the rational Betti vector $F\in\operatorname{End}^0(H_B)$ is
\[
 S=P_c^{-1}FP_c
 =\frac1D\begin{pmatrix}xq+yp&2pq\\-2xy&-xq-yp\end{pmatrix}.
\]
For $a\in\Qb$ and the algebraic functional $\ell_a(T)=T_{11}-aT_{21}$,
\[
 h=a\quad\Longleftrightarrow\quad
 \ell_a(S)=\frac{xq+yp+2axy}{D}=0.
\]
This is a period relation on
$\operatorname{End}^0(H)\simeq\operatorname{Sym}^2H(1)$.
\end{result}
\begin{proof}
The columns $(1,\tau)$ and $(1,\bar\tau)$ are the two Hodge lines in Betti coordinates. Conjugating $\operatorname{diag}(0,1)$ by their column matrix gives the asserted formula. Conversely,
$\tau=-(e_B)_{11}/(e_B)_{12}$ and $\bar\tau=(e_B)_{22}/(e_B)_{12}$, proving the field equality. Schneider's theorem, already applied in A.3.1, proves the CM assertion. Direct matrix multiplication gives $S$ and $\ell_a(S)$; the ratio $S_{11}/S_{21}=h$ was established above. Polarization gives $H^\vee\simeq H(1)$, and
$H\otimes H=\operatorname{Sym}^2H\oplus\bigwedge^2H$, with $\bigwedge^2H(1)=\Q(0)$; the trace-free summand is therefore the stated rank-three tensor.
\end{proof}

Algebraicity of $h$ verifies that $e_{\dR}$ is algebraic. Formula A.4.11 computes the additional Betti coefficient-field requirement exactly. Similarly, the equation $\ell_a(S)=0$ places $S$ in an algebraic de Rham plane, whereas a de Rham--Betti class represented by the rational vector $F$ would require all coordinates of $S$ to be algebraic. These are explicit checks of the candidate tensors. They make no assertion that no other construction could help.

\begin{result}{A.4.12}{The actual algebraic-splitting hypersurface generates $\mathrm{GL}_2$}{Proved}
For $a\in\Qb$, the variety
\[
 V_a=\left\{\begin{pmatrix}x&p\\y&q\end{pmatrix}\in\mathrm{GL}_2:
 xq+yp+2axy=0\right\}
\]
is irreducible of dimension three. The smallest algebraic subtorsor of $\mathrm{GL}_2$ containing $V_a$ is $\mathrm{GL}_2$.
\end{result}
\begin{proof}
Put $B_a=\left(\begin{smallmatrix}1&a\\0&1\end{smallmatrix}\right)$, so that $V_aB_a=V_0$, and put $H=\left(\begin{smallmatrix}1&1\\1&-1\end{smallmatrix}\right)$. Every element of $V_0$ has $x,y,p\ne0$ and admits the unique parametrization
\[
 \begin{pmatrix}x&p\\y&-yp/x\end{pmatrix}
 =\operatorname{diag}(x,y)H\operatorname{diag}(1,p/x).
\]
Its determinant is $-2yp$, giving $V_0\simeq\Gm^3$. Let $T$ be the diagonal torus. The translated set $H^{-1}V_0$ contains $T$ (take $x=y$) and $H^{-1}TH$ (take $p=x$). The Lie algebra of its generated closed subgroup contains
\[
 I,\quad Z=\begin{pmatrix}1&0\\0&-1\end{pmatrix},\quad
 X=H^{-1}ZH=\begin{pmatrix}0&1\\1&0\end{pmatrix},\quad
 [Z,X]=\begin{pmatrix}0&2\\-2&0\end{pmatrix}.
\]
They form a basis of $\mathfrak{gl}_2$, so the closed generated subgroup has dimension four and equals the connected group $\mathrm{GL}_2$. Translation by $B_a$ and the corresponding change of basepoint give the claim for $V_a$.
\end{proof}

Kreutz--Shen--Vial \cite[Proposition 3.1; Theorem 5.13(ii)]{GenKSV2026}, with Betti coefficient field $L=\Qb$, yield the full generated comparison torsor for one non-CM elliptic curve, as proved in A.4(b). Result A.4.12 shows that containment in the precise hypersurface arising from $h=a$ does not contradict that conclusion. No assertion that $V_a$ is an actual elliptic period locus is made. The symbolic checker \texttt{code/check\_torsor\_gate.py} verifies the matrices, factorization, Lie-algebra rank, and projector identities exactly.

\begin{result}{A.4.13}{Exact applicability of the functional theorem to a single splitting value}{Proved}
In the Legendre family
\[
 E_t:Y^2=4X(X-1)(X-t),\qquad t\in S=\A^1\setminus\{0,1\},
\]
use $x=X-(1+t)/3$, $\omega=dx/Y$, and $\xi=x\,dx/Y$. Choose period cycles near a real $t_0\in(0,1)$ with first-row periods real and second-row periods purely imaginary along the real interval. For any fixed $a\in\Qb$, the relation $h(t_0)=a$ gives an isolated intersection of the horizontal period leaf with
\[
 W_a=\{R_a(P)=0\}\subset\Omega,\qquad
 R_a\begin{pmatrix}x&p\\y&q\end{pmatrix}=xq+yp+2axy.
\]
Here $\dim\Omega=4$, $\dim W_a=3$, and the algebraic monodromy group has dimension three. Thus this isolated intersection has codimension equal to the monodromy dimension, and does not satisfy the strict inequality in Bakker--Tsimerman Theorem~1.3. Adding $t=t_0$ to $W_a$ makes the strict inequality hold, but its conclusion is satisfied by every point $t_0$, without requiring CM.
\end{result}
\begin{proof}
The displayed translation gives
\[
 g_2(t)=\tfrac43(t^2-t+1),\qquad
 g_3(t)=\tfrac4{27}(2-3t-3t^2+2t^3),
\]
so A.4.11 supplies the relation $R_a=0$ on the real interval. The periods are proportional to $K(t)$ and $iK(1-t)$. The logarithmic term of $K(1-t)$ at zero is $-K(t)\log(t)/\pi$: the coefficient follows from its Picard--Fuchs equation and $K(1-t)\sim\tfrac12\log(16/t)$. Continuation around zero adds twice the first period to the second; the corresponding calculation at one gives the other unipotent. The two one-parameter unipotent groups in their Zariski closure have Lie algebras whose bracket spans $\mathfrak{sl}_2$. Monodromy is therefore Zariski-dense in $\mathrm{SL}_2$.

The period torsor has determinant-fixed three-dimensional fibers and one-dimensional base. Put $D=2\pi i$. Its hypersurface $W_a$ has the fiber parametrization
\[
 P=\begin{pmatrix}
 x&-D/(2y)-ax\\y&D/(2x)-ay
 \end{pmatrix},\qquad x,y\ne0,
\]
so its total dimension is three. The function $R_a(P(t))$ is holomorphic on a local period branch and cannot vanish identically: otherwise continuation around loops, followed by Zariski density of monodromy, would put a whole determinant fiber in the proper two-dimensional fiber just parametrized. Every zero is therefore isolated. Its codimension in $W_a$ is three, equal to $\dim\mathrm{SL}_2$; Theorem~1.3 of \cite{EllBakkerTsimerman2025}, p.5, assumes a strict inequality.

Imposing $t=t_0$ gives a two-dimensional $W$. The theorem then places its projected point in a proper weak Mumford--Tate subvariety. Corollary~2.9 of the same source, p.8, characterizes that containment by a drop in algebraic monodromy. Restriction to any point has trivial monodromy, strictly smaller than $\mathrm{SL}_2$. Hence this conclusion applies to every point and supplies no CM condition.
\end{proof}

In the formulation of \cite[Theorem 4.25, p.18]{EllBakkerTsimerman2025}, spreading a fixed elliptic curve constantly over $\C(t)$ gives trivial relative monodromy and relative motivic group. Its period functions already belong to the constant field $\C$, so the lower bound is zero. Spreading in the nonconstant Legendre family instead gives the preceding calculation; an equality at one fiber does not become an identity of meromorphic functions. These conclusions concern the specified applications. The gate for a nonzero algebraic $h$ remains unresolved.

\begin{result}{A.4.14}{The exact unitary lift in the universal vector extension}{Proved}
Let $c=\pi/A$. The universal vector extension has the uniformization
\[
 E^\dagger(\C)=\C^2/\Lambda^\dagger,\qquad
 \Lambda^\dagger=\Z(\omega_1,\eta_1)+\Z(\omega_2,\eta_2).
\]
On its algebraic chart above $E\setminus\{O\}$, the coordinates and invariant differentials are
\[
 (u,s)\longmapsto(\wp(u),\wp'(u),t=\zeta(u)-s),\qquad
 du=dx/y,\quad ds=-dt-x\,dx/y.
\]
The canonical real Lie-group section $\kappa:E(\C)\to E^\dagger(\C)$, with image the maximal compact subgroup, is
\[
 \kappa([u])=[(u,hu+c\bar u)].
\]
For every $P=[u]\in E(\Qb)\setminus\{O\}$,
\[
 t(\kappa(P))=\zeta(u)-hu-c\bar u\in\Qb
 \quad\Longleftrightarrow\quad P\text{ is torsion}.
\]
The complex-linear part of $d\kappa$ has image $L_h=\{s=hu\}\subset\operatorname{Lie}E^\dagger$. This line is defined over $\Qb$ exactly when $h$ is algebraic. Its annihilator $ds-hdu$ represents $-(\xi+h\omega)$ in de Rham cohomology.
\end{result}
\begin{proof}
The algebraic chart and differentials are \cite[Sections 4.1--4.2, equation (4.1) and Proposition 4.2, pp.7--8]{EllFonsecaMatthes2020}. Their period of $x\,dx/y$ has the negative of our quasiperiod sign, giving precisely the lattice displayed here. By A.4.1,
$\eta_i=h\omega_i+c\bar\omega_i$. Consequently the real span of $\Lambda^\dagger$ is
\[
 W=\{(u,hu+c\bar u):u\in\C\}.
\]
The quotient $W/\Lambda^\dagger$ is compact. Every compact subgroup of $\C^2/\Lambda^\dagger$ maps trivially into the real vector group $\C^2/W$, so this compact subgroup is maximal. Projection to $\C/\Lambda$ is an isomorphism of real Lie groups, proving the formula for $\kappa$. The displayed $t$-coordinate is unchanged when $u$ is replaced by $u+\omega_i$, so it is well-defined.

Bost--K\"unnemann \cite[Theorem 3.1.2, p.386]{EllBostKunnemann2009} prove that an algebraic point of a universal vector extension which belongs to its maximal compact subgroup at one complex embedding is torsion. Their Appendix B.7, equations (B.17)--(B.20), pp.421--422, identifies this subgroup with line bundles equipped with their unitary connection. If the displayed $t$-value is algebraic, all three coordinates of $\kappa(P)$ are algebraic and the theorem makes $P$ torsion. Conversely, multiplication by $n$ is invertible on the additive kernel, so $E^\dagger[n]\to E[n]$ is an isomorphism. The unique torsion lift of a torsion $P$ is algebraic and lies in the maximal compact subgroup; hence it is $\kappa(P)$.

Differentiating $s=hu+c\bar u$ gives complex-linear part $ds=hdu$, and the two differentials $du,ds$ are algebraic. Their coordinate expression gives $[ds-hdu]=-[\xi+h\omega]$. This proves the last assertions.
\end{proof}

The pointwise arithmetic hypothesis of this theorem is the algebraicity of the entire value $\zeta(u)-hu-c\bar u$. Algebraicity of $h$ alone specifies the tangent line $L_h$. For this same line, $\exp(L_h)$ intersects the maximal compact subgroup only at the identity: a vector $(u,hu)$ belongs to $W$ precisely when $c\bar u=0$, hence $u=0$, and adding a lattice vector preserves membership in $W$. These formulas identify the exact additional quantity needed to use the unitary-connection theorem at a non-torsion algebraic point.

\begin{result}{A.4.15}{The restricted-Lie test for an algebraic splitting line}{Proved}
For every $a\in\Qb$, the line $L_a=\{s=au\}$ in $\operatorname{Lie}E^\dagger$ is not the Lie algebra of an algebraic subgroup. Its reductions therefore fail to be closed under $p$-th powers at infinitely many good primes of every number field of definition.

In the preceding algebraic chart, put
\[
 D_a=y\partial_x+(6x^2-g_2/2)\partial_y-(x+a)\partial_t,
 \qquad V=-\partial_t.
\]
The derivation $D_a$ generates $L_a$. In characteristic five,
\[
 D_a^5=2g_2D_a+(3g_3-2g_2a)V.
\]
In particular, the actual CM example of A.4.4, with $(g_2,g_3,h)=(30,28,1/2)$, has good reduction at five and satisfies $D_h^5=4V$, which is not in the line generated by $D_h$.
\end{result}
\begin{proof}
A connected one-dimensional algebraic subgroup with nonzero projection to $E$ would map surjectively to $E$, with finite kernel contained in $\Ga$. In characteristic zero this kernel is trivial, so the subgroup would split $E^\dagger\to E$. A holomorphic splitting would lift to a complex-linear map $u\mapsto(u,bu)$ on universal covers and would require $b\omega_i=\eta_i$ for both $i$. This contradicts $\eta_1\omega_2-\eta_2\omega_1=2\pi i$. The only remaining one-dimensional algebraic subgroup is vertical, so its Lie algebra is not $L_a$.

Bost's restricted-Lie criterion \cite[Theorem 2.3, p.172]{EllBost2001} states that a Lie subalgebra of an algebraic group over a number field is algebraic exactly when its reductions are closed under $p$-th powers at all but finitely many primes. Its contrapositive gives the infinite-prime assertion.

The invariant forms give $du(D_a)=1$ and $ds(D_a)=a$. In characteristic five both sides of the claimed identity are derivations of the coordinate ring $y^2=4x^3-g_2x-g_3$. Checking its three generators suffices. Their values on $x,y,t$, respectively, are
\[
 2g_2y,\qquad 2g_2x^2+4g_2^2,\qquad 3g_2x+2g_3.
\]
This proves the identity. For the CM example the discriminant is $5832$, a unit modulo five, $g_2=0$ and $g_3=3$ modulo five, so $D_h^5=4V$. Since $du(D_h)=1$ but $du(V)=0$, the latter is not a multiple of $D_h$. The computation is also checked exactly in \texttt{code/check\_p\_curvature.py}.
\end{proof}

Thus algebraicity of the Hodge complement does not supply the almost-everywhere restricted-Lie hypothesis for this particular line: A.4.15 proves its failure even for the proved CM examples. This tests the specified algebraization attempt; it does not exclude other uses of arithmetic foliations or $p$-adic comparison maps. The nonzero-value gate for $h$ remains unresolved.

\begin{result}{A.4.16}{Torsion lifts and the exact horizontal-lift test}{Proved}
Write $Z(P)=\zeta(u)-hu-c\bar u$ for $P=[u]\ne O$. The addition and multiplication formulas are
\[
 Z(P+Q)=Z(P)+Z(Q)+\frac{y(P)-y(Q)}{2(x(P)-x(Q))},
 \qquad Z(mP)=mZ(P)+R_m(P),
\]
where $R_m\in\Q(g_2,g_3)(E)$ is independent of $h$ and $c$, and the formulas hold wherever defined. For $P$ of exact order $n\ge2$,
\[
 Z(P)=-R_{n-1}(P)/n.
\]
In particular, $Z(P)=0$ at nonzero two-torsion, and
\[
 Z(P)=-\frac{12x(P)^2-g_2}{12y(P)}
\]
at nonzero three-torsion. For every algebraic $a$ and every nonzero algebraic point $P=[u]$, however,
\[
 \exp(L_a)\cap E^\dagger(\Qb)=\{O\},\qquad
 \zeta(u)-au\notin\Qb.
\]
\end{result}
\begin{proof}
The zeta addition formula gives the first identity because the two real-linear correction terms cancel. Its duplication limit is
\[
 Z(2P)=2Z(P)+\frac{12x(P)^2-g_2}{4y(P)}.
\]
Induction using these formulas expresses $R_m(P)=\zeta(mu)-m\zeta(u)$ as a rational function over $\Q(g_2,g_3)$. Equivalently, its quasiperiod increments cancel, so it is elliptic; the addition recursion verifies its coefficient field. Since $(n-1)P=-P$ and $Z$ is odd, the torsion formula follows. Taking $n=2,3$ gives the displayed special values.

For the second assertion, apply W\"ustholz's analytic subgroup theorem to
\[
 \Ga(\C)\longrightarrow E^\dagger(\C),\qquad
 u\longmapsto\exp(u,au).
\]
Its differential is algebraic by A.4.14. The original primary preprint \cite[Hauptsatz and Bemerkung 2, printed p.3]{EllWustholz1984} states that a nonidentity algebraic point in the image forces a positive-dimensional algebraic subgroup in the image. Such a subgroup would have Lie algebra $L_a$, contradicting A.4.15. Thus the only algebraic point in the image is the identity. Above $P\ne O$ its chart coordinates are $(x(P),y(P),\zeta(u)-au)$, proving the last assertion.
\end{proof}

\begin{result}{A.4.17}{An algebraic compact point on the explicit non-CM curve}{Proved}
Take $E_*:y^2=4x^3-4x-1$. Let $r\in(1,2)$ be the unique root there of $3r^4-6r^2-3r-1=0$, and put $s=\sqrt{4r^3-4r-1}>0$. Then
\[
 P=(r,s),\qquad
 \kappa(P)=\left(r,s,-\frac{3r^2-1}{3s}\right)
\]
are algebraic points, $P$ has exact order three, and $E_*$ is non-CM.
\end{result}
\begin{proof}
The duplication slope $m=(3r^2-1)/s$ satisfies $m^2=3r$, so $2P=-P$; $s\ne0$ follows from the resultant $-1369$ of the two defining polynomials. Non-CM follows from $j(E_*)=110592/37$. The lift's last coordinate $t$ satisfies
\[
 243t^8-54t^4-9t^2-1=0.
\]
These identities are checked in \texttt{code/check\_torsion\_lifts.py}.
\end{proof}

The addition and division operations starting from torsion points produce only torsion points: a finite sum has a common annihilating integer, and $nQ=P$, $mP=0$ imply $mnQ=0$. Their completed-zeta recurrences therefore give algebraic compact points for every algebraic elliptic curve without supplying a non-torsion point. The alternative horizontal lift is tested by the last assertion of A.4.16. These are exact calculations for the proposed constructions; the general non-CM splitting gate is unchanged.

\begin{result}{A.4.18}{The exact de Rham Hodge endomorphism algebra}{Proved}
Define
\[
 \mathcal D_E=\{T\in\operatorname{End}_{\Qb}H^1_{\dR}(E/\Qb):
 T_\C H^{1,0}\subset H^{1,0},\quad T_\C H^{0,1}\subset H^{0,1}\}.
\]
Writing $e=e_{\dR}$ as in A.4.11 gives the exact formula
\[
 \mathcal D_E
 =\{aI+(d-a)e:a,d\in\Qb,\ (d-a)h\in\Qb\}.
\]
Consequently $\mathcal D_E=\Qb I\oplus\Qb e\simeq\Qb\times\Qb$ when $h\in\Qb$, and $\mathcal D_E=\Qb I$ otherwise. The injective pullback map
\[
 \iota:\operatorname{End}_{\Qb}(E)\otimes_{\Z}\Qb
 \longrightarrow\mathcal D_E
\]
has precisely the image
\[
 \iota\bigl(\operatorname{End}_{\Qb}(E)\otimes\Qb\bigr)
 =\{T\in\mathcal D_E:T_B\in M_2(\Qb)\}.
\]
For CM curves this map is surjective. For a non-CM curve its image is $\Qb I$, so its cokernel has dimension one exactly when $h$ is algebraic. The Hodge projector $e_B$ itself is never rational over $\Q$, including in the CM case.
\end{result}
\begin{proof}
Preservation of $H^{1,0}=\C\omega$ first gives
$T=\left(\begin{smallmatrix}a&b\\0&d\end{smallmatrix}\right)$ with algebraic entries. Preservation of $\C(h\omega+\xi)$ gives $ah+b=dh$, or $b=(d-a)h$. This proves the formula and both descriptions of $\mathcal D_E$.

Pullback by an algebraic endomorphism preserves both Hodge lines and has rational Betti matrix; after tensoring, its Betti matrix has algebraic entries. Faithfulness follows already from the action on the complex torus: an endomorphism acting trivially on $H^1$ has zero differential and is zero. In the CM case the endomorphism field has degree two over $\Q$, and its two embeddings act with distinct eigenvalues on the two Hodge lines. Its scalar extension to $\Qb$ therefore maps onto the two-dimensional algebra $\mathcal D_E$.

For the opposite inclusion, take $T=aI+(d-a)e\in\mathcal D_E$ with algebraic Betti entries. If $d=a$, then $T=aI$ lies in the image. If $d\ne a$, then
\[
 e_B=\frac{T_B-aI}{d-a}\in M_2(\Qb).
\]
By A.4.11 this forces CM, and the preceding CM calculation proves membership in the image. For non-CM curves $\operatorname{End}_{\Qb}(E)=\Z$, giving the asserted cokernel description. Finally,
\[
 (e_B)_{12}=\frac1{\bar\tau-\tau}=\frac{i}{2\operatorname{Im}\tau}
\]
is nonzero and purely imaginary, so $e_B$ cannot have rational entries.
\end{proof}

Bost--Charles \cite[Theorem 3.8(1), equation (3.11), p.17]{EllBostCharles2016} prove full faithfulness for the de Rham--Betti realization with its integral Betti lattice; after inverting integers its endomorphisms have rational Betti matrices. The map just calculated identifies the additional coefficient-field condition exactly, including its version with algebraic Betti coefficients. Deligne's absolute-Hodge theorem \cite[Section 2, definitions on pp.15--16 and Main Theorem 2.11, p.19]{EllDeligne1982} starts with a rational Betti Hodge class. The nonzero imaginary entry above proves that this hypothesis fails for the single projector $e$, even when $h$ is algebraic on a CM curve. On that curve $e$ belongs to the algebraic-coefficient span of genuine endomorphisms. The scalar extension from $\Q$ to $\Qb$ is therefore essential. No converse about the algebraicity of a non-CM value of $h$ is supplied by this calculation.

\begin{result}{A.4.19}{Crystalline tests of the same projector, with coefficient fields retained}{Proved}
Let $E$ and a rank-one projector $e\in\operatorname{End}_L H^1_{\dR}(E/L)$ be defined over a number field $L$. Its crystalline-Frobenius fixity set cannot have lower weighted density greater than $3/4$, where the density is
\[
 \underline\delta(e)=\liminf_{X\to\infty}
 \frac{\displaystyle\sum_{\substack{v:\,p_v\le X\\\varphi_v(e)=e}}
 [L_v:\Q_{p_v}]\,\frac{\log p_v}{p_v-1}}
 {[L:\Q]\displaystyle\sum_{p\le X}\frac{\log p}{p-1}}.
\]
Finitely many ramified or bad-reduction places are omitted. In particular such a projector is never a $1$-absolute-Tate endomorphism.

At an unramified good place, write $\varphi_v=F_v\sigma_v$, where $\sigma_v$ acts on the coefficients and $F_v=\left(\begin{smallmatrix}A&B\\C&D\end{smallmatrix}\right)$. For $e=e_h=\left(\begin{smallmatrix}0&h\\0&1\end{smallmatrix}\right)$ defined over $L$, the exact fixity test is
\[
 \varphi_v(e_h)=e_h
 \quad\Longleftrightarrow\quad
 C=0\quad\hbox{and}\quad A\sigma_v(h)+B=hD.
\]
For the CM curve $E:y^2=4x^3-30x-28$ and its $\Q$-defined projector $e_{1/2}$, this test fails at $p=5$.
\end{result}
\begin{proof}
Tang \cite[Definition 2.3.1 and Theorem 2.3.3, p.7]{EllTang2017} proves that the $1$-absolute-Tate endomorphisms over $L$ are precisely the image of
$\operatorname{End}_L(E)\otimes_{\Z}\Q$. Corollary~5.2.3 of the same paper, p.23, strengthens this: for a geometrically simple abelian variety, an element of $\operatorname{End}_L H^1_{\dR}(E/L)$ with fixity density greater than $3/4$ is an algebraic endomorphism with rational coefficients. Every elliptic curve is geometrically simple. Every nonzero element of $\operatorname{End}_{\Qb}(E)\otimes\Q$ is invertible, because a nonzero elliptic-curve endomorphism is an isogeny. This rational algebra therefore has no nontrivial idempotent. A rank-one projector is a nontrivial idempotent, giving the density bound.

This argument uses $\operatorname{End}(E)\otimes\Q$, not $\operatorname{End}(E)\otimes\Qb$. In the CM case the latter is a product of two copies of $\Qb$ and does contain the Hodge projectors of A.4.18. Crystalline Frobenius acts semilinearly on their algebraic coefficients. Indeed, direct multiplication gives
\[
 F_v\sigma_v(e_h)=
 \begin{pmatrix}0&A\sigma_v(h)+B\\0&C\sigma_v(h)+D\end{pmatrix},
 \qquad
 e_hF_v=\begin{pmatrix}hC&hD\\C&D\end{pmatrix}.
\]
Equality is exactly the displayed pair of conditions.

For the specified curve over $\Q$, the discriminant $5832$ is a unit at five. Its reduction is $y^2=4x^3+2$ over $\mathbf F_5$. For $x=0,1,2,3,4$, the numbers of solutions in $y$ are respectively $0,2,2,1,0$. Including the point at infinity gives six points, so the Frobenius trace is zero and its characteristic polynomial on crystalline $H^1$ over $\Q_5$ is $T^2+5$. This polynomial is irreducible over $\Q_5$, since a square has even valuation whereas $-5$ has odd valuation. There is no invariant $\Q_5$-line. A Frobenius-fixed $\Q_5$-linear projector would make its image such a line, a contradiction. Here the base field is $\Q$ and $\sigma_5$ is the identity. No replacement by linear Frobenius after extending to a CM coefficient field is made.
\end{proof}

For the actual CM endomorphism of A.4.4, let $u^2=-2$. Its de Rham matrix is
$\alpha_{\dR}=u\left(\begin{smallmatrix}1&-1\\0&-1\end{smallmatrix}\right)$ and
$e_{1/2}=(\alpha_{\dR}-uI)/(-2u)$. This displays the algebraic coefficients that create the projector from a genuine endomorphism. The preceding Frobenius calculation tests the proposed arithmetic-Hodge construction explicitly, while leaving the nonzero non-CM splitting gate unresolved.

\begin{result}{A.4.20}{The Frobenius condition at one good place}{Proved}
Let $E/L$ be a non-CM elliptic curve over a number field. At every unramified good-reduction place $v$ of residue characteristic $p>2$, absolute crystalline Frobenius does not preserve
\[
 \operatorname{Fil}^1 H^1_{\dR}(E/L_v).
\]
Consequently no rank-one de Rham projector whose kernel is this Hodge line can be fixed by that Frobenius. This applies to the projector of A.4.11 whenever its coefficients belong to $L$.
\end{result}
\begin{proof}
Write $k_v=\mathbf F_q$, $q=p^f$, and $W=W(k_v)=\mathcal O_{L_v}$; the equality uses the unramified hypothesis. Put
\[
 M=H^1_{\dR}(\mathcal E/W)=H^1_{\mathrm{cris}}(E_{k_v}/W),
 \qquad F=\varphi_v.
\]
The Hodge line is a saturated direct summand of the rank-two $W$-module $M$. Suppose its fraction-field extension is $F$-stable. Then the integral line is stable as well, since $F(M)\subset M$. In a basis beginning with its generator, the semilinear Frobenius has upper-triangular matrix
\[
 \begin{pmatrix}a&b\\0&d\end{pmatrix},\qquad a,b,d\in W.
\]
Modulo $p$, Frobenius pullback annihilates invariant differentials, so $a\in pW$. Cup product identifies $\bigwedge^2M$ with the rank-one crystalline $H^2$, on which Frobenius acts as $p\sigma_v$. Thus $v_p(ad)=1$, and integrality gives $v_p(a)=1$, $v_p(d)=0$. The linear iterate $F^f$ has diagonal entries
\[
 \prod_{j=0}^{f-1}\sigma_v^j(a),\qquad
 \prod_{j=0}^{f-1}\sigma_v^j(d),
\]
of valuations $f$ and zero. It is therefore nonscalar.

The endomorphism $\pi_q$ of the special fiber has crystalline pullback $F^f$, which preserves the integral Hodge filtration. The Serre--Tate--Grothendieck--Messing lifting criterion therefore lifts $\pi_q$ to an endomorphism of $\mathcal E/W$. Here the divided powers on the ideal $pW/p^n$ are nilpotent: for $p>2$,
\[
v_p(p^m/m!)=m-v_p(m!)\geq m\frac{p-2}{p-1}\longrightarrow\infty.
\]
Grothendieck--Messing first lifts the $p$-divisible-group morphism over each $W/p^n$; the Serre--Tate equivalence lifts the actual special-fiber abelian endomorphism uniquely. Uniqueness makes these lifts compatible, and proper formal existence algebraizes their graphs to an endomorphism over $W$. The exact functor and its morphism condition are stated and proved in \cite[Theorem 3.4 and Remark 3.5, pp.12--13, and Appendix B]{EllPannyRobertSferlazza2026}; the original Serre--Tate fully faithful deformation functor and ordinary homomorphism criterion are \cite[Theorem 1.2.1 and Theorem 2.1(4), pp.143--145 and 149]{EllKatz1981}. The lifted endomorphism is not an integer multiplication, since its crystalline action is nonscalar.

For completeness, this local endomorphism gives an endomorphism over $\Qb$. Its finite kernel consists of torsion points of the original algebraic elliptic curve, hence is defined over $\Qb$ inside $\overline{L_v}$. The quotient by that kernel is defined over $\Qb$. Its isomorphism to $E$ also descends: equal $j$-invariants give an isomorphism over the algebraically closed field $\Qb$, and the difference from the given isomorphism is an automorphism of $E$, all of which are defined over $\Qb$. A noninteger endomorphism in characteristic zero gives CM, contrary to the hypothesis. This proves failure of Hodge-line preservation. A Frobenius-fixed projector preserves its kernel under Frobenius, so the last assertion follows.
\end{proof}

This result supplies the exact endomorphism produced by a Frobenius-stability argument and a strict criterion that would have to be established at a place before that argument applies. Algebraicity of the complex number $h$ has not been shown to impose this additional crystalline equation. The CM example at five in A.4.19 proves that it cannot impose it at every unramified good place.

\begin{result}{A.4.21}{An explicit ordinary Hodge--Tate base change}{Proved}
For $E_*:y^2=4x^3-4x-1$, let
\[
 G=\mathcal E_*[5^\infty]\otimes_{\Z_5}\mathcal O_{\C_5}.
\]
Then $G$ has CM in the sense of one-dimensional $p$-divisible groups, and its Hodge--Tate period field is $\Q_5$. The elliptic curve $E_*\otimes\C_5$ has no CM.
\end{result}
\begin{proof}
The discriminant is $37$. At five, the numbers of $y$-solutions at $x=0,1,2,3,4$ are $2,2,0,1,2$. Thus $\#E_*(\mathbf F_5)=8$ and $a_5=-2$, giving ordinary good reduction. Over $\mathcal O_{\C_5}$ its connected--\'etale sequence has height-one factors and is
\[
 0\longrightarrow\mu_{5^\infty}\longrightarrow G
 \longrightarrow\Q_5/\Z_5\longrightarrow0.
\]
Here the unramified twists of the two height-one factors have become trivial after this base change. In Serre--Tate coordinates the extension is specified by a unit $q_{\mathrm{ST}}\in1+\mathfrak m_{\C_5}$. Choose compatible roots $q_n$ with $q_n^{5^n}=q_{\mathrm{ST}}$ and $q_{n+1}^5=q_n$. In the Kummer presentation of $G[5^n]$, the component indexed by $a\bmod5^n$ has equation $x^{5^n}=q_{\mathrm{ST}}^a$. The points $(q_n^a,a)$ form a section: addition of components uses the carry factor $q_{\mathrm{ST}}^{-\lfloor(a+b)/5^n\rfloor}$, exactly cancelled by $q_n^{5^n}=q_{\mathrm{ST}}$. The inclusion into level $5^{n+1}$ sends $(x,a)$ to $(x,5a)$, so the chosen sections are compatible. Consequently
\[
 G\simeq\mu_{5^\infty}\oplus\Q_5/\Z_5.
\]
The two projections give the rank-two semisimple algebra $\Q_5\times\Q_5$ inside $\operatorname{End}(G)\otimes\Q_5$, establishing CM by the definition in \cite[p.673]{EllHowe2018}. The exact Hodge--Tate sequence is
\[
 0\longrightarrow\operatorname{Lie}G(1)\otimes\C_5
 \longrightarrow T_5G\otimes\C_5
 \longrightarrow\omega_{G^\vee}\otimes\C_5\longrightarrow0.
\]
In a Tate-module basis of the two displayed summands, its Hodge--Tate line is the multiplicative summand, with coordinate $[1:0]$. Its field of definition is therefore $\Q_5$. This is the sequence in \cite[(1.1), p.672]{EllHowe2018}; Theorem~1.2 and Remark~1.3 on p.673 explicitly permit this ordinary splitting after base change.

On the other hand, $j(E_*)=110592/37$ is not an algebraic integer, so $E_*$ is non-CM. Its endomorphism algebra is unchanged by extension to an algebraically closed field, by the finite-kernel and quotient descent argument in A.4.20. Thus the base-changed elliptic curve still has no CM.
\end{proof}

The functor in this example is the actual map $E\mapsto E[5^\infty]\otimes\mathcal O_{\C_5}$, and the extra endomorphisms are the explicitly constructed projections after its ordinary extension splits. Howe's theorem correctly detects CM of this $p$-divisible group; the same explicit curve proves that this conclusion alone does not yield an elliptic-curve endomorphism. Its complex splitting coefficient $h$ remains subject to the original unresolved gate.

\begin{result}{A.4.22}{The exact denominator condition for an ordinary correction}{Proved}
Let $f(q)=\sum_{n\ge1}b(n)q^n$ be a normalized rational weight-two newform, let
\[
 \varepsilon(q)=\sum_{n\ge1}\frac{b(n)}n q^n,
 \qquad Z_a(q)=\zeta(\Lambda_f,\varepsilon(q))-a\varepsilon(q),
\]
and use the elliptic period lattice $\Lambda_f$ of Guerzhoy's construction. For every good ordinary prime outside its finite exceptional set there is $\lambda_p\in\Z_p$ such that, for every $a\in\C_p$,
\[
 \inf_n v_p([q^n]Z_a)>-\infty
 \quad\Longleftrightarrow\quad a=\lambda_p.
\]
Algebraic coefficients do not, even in an actual CM splitting example, ensure bounded denominators. On
\[
 E_0:y^2=4x^3-4x,\qquad h(E_0)=0,
\]
put $t=-2x/y$ and let $u(t)$ be the formal logarithm with $du=dx/y$. Then
\[
 \zeta(u(t))=\frac1t-\frac12\sum_{n\ge1}
 \frac{(-1)^n\binom{2n}{n}}{4n-1}t^{4n-1}\in\Q((t)),
\]
and its coefficients have unbounded denominators at three. More precisely, for
$n_r=(3^{2r+1}+1)/4$, the coefficient of $t^{3^{2r+1}}$ has valuation $-r-1$.
\end{result}
\begin{proof}
Guerzhoy \cite[Theorem 2, pp.4--5, equation (4)]{EllGuerzhoy2014} constructs
\[
 \zeta(\Lambda_f,\varepsilon)-\lambda_p\varepsilon
 -\mu_p\varepsilon(q^p)\in\Z_p((q)),
\]
with $\lambda_p\in\Z_p$ and $\mu_p=0$ at the ordinary primes under consideration. The inequality in that last hypothesis is $b(p)\not\equiv0\pmod p$, as verified on the original PDF page; text extraction can omit its slash. Thus $Z_{\lambda_p}$ is integral. The Hecke recurrence
\[
 b(p^r)=b(p)b(p^{r-1})-p\,b(p^{r-2})
\]
gives $b(p^r)\equiv b(p)^r\pmod p$, so
$v_p([q^{p^r}]\varepsilon)=-r$. For $a\ne\lambda_p$, the coefficient at $q^{p^r}$ in
\[
 Z_a=Z_{\lambda_p}+(\lambda_p-a)\varepsilon
\]
therefore has valuation $v_p(\lambda_p-a)-r$ for all sufficiently large $r$: it is then strictly smaller than the valuation of the integral first summand. These valuations tend to minus infinity. This proves both directions of the equivalence.

For the CM example, A.4.6 gives $h(E_0)=0$. Write $R=\sqrt{1+4t^4}$, with $R(0)=1$. The curve equation and $t=-2x/y$ give
\[
 x=\frac{1+R}{2t^2},\qquad
 \frac{dx}{y}=\frac{dt}{R},\qquad
 x\frac{dx}{y}=\frac{1+R^{-1}}{2t^2}\,dt.
\]
Since $d\zeta=-x\,du$ and $\zeta(u(t))=t^{-1}+O(t)$, integration of
$R^{-1}=\sum_{n\ge0}(-1)^n\binom{2n}{n}t^{4n}$ gives the asserted series. There is no integration constant, since both $u(t)$ and $\zeta$ are odd.

For $r\ge0$,
\[
 n_r=1+2\sum_{j=0}^{r-1}3^{2j+1}.
\]
Adding $n_r$ to itself in base three creates exactly $r$ carries: one from each odd-position digit two, and none from the intervening zero digits. The elementary factorial-valuation formula, equivalently its carry form, gives
$v_3\binom{2n_r}{n_r}=r$. As $4n_r-1=3^{2r+1}$, the displayed coefficient has valuation
$r-(2r+1)=-r-1$. This proves unboundedness on the stated algebraic CM curve. The local parameter is integral and has a unit linear term at three; integral invertible changes of formal parameter preserve boundedness, by applying substitution and its integral inverse coefficientwise.
\end{proof}

For a putative algebraic value of the complex $h$, an embedding of its number field into $\C_p$ makes the first displayed criterion directly applicable. The additional equation to establish is $\iota_p(h)=\lambda_p$, equivalently bounded denominators of that particular embedded series. The proof has not established this equation. Guerzhoy's adele records its separate complex and finite components; membership in the adele ring does not identify them with a single diagonal algebraic number. The supersingular CM computation proves that an unrestricted passage from algebraic coefficients to bounded denominators would be false.

\begin{result}{A.4.23}{The geometric mock-form coefficient is the same splitting scalar}{Proved}
Let $\pi:X_0(N)\to E$ be a modular parametrization over $\Q$, of degree $d$, with
\[
 \pi^*\omega=m\omega_f,\qquad
 \omega_f=f(q)\frac{dq}{q},\qquad m\in\Q^\times,
\]
for a normalized rational newform $f$. In its two-dimensional de Rham component choose
\[
 \eta_f^{\rm alg}=\frac m d\pi^*\xi.
\]
This class is cup-dual to $\omega_f$. The complex number $a$ defined by Candelori's decomposition
\[
 \eta_f^{\rm alg}=a\omega_f+b\overline{\omega_f}
\]
satisfies the exact identity
\[
 a=-\frac{m^2}{d}h(E).
\]
For the weight-zero harmonic Maass form constructed from this class, with algebraic constant term, all of its holomorphic coefficients are algebraic exactly when $h(E)$ is algebraic. With a rational constant term, their field is precisely $\Q(h(E))$.
\end{result}
\begin{proof}
Our period convention gives $\langle\omega,\xi\rangle=1$: the determinant of their periods is $2\pi i$, and the de Rham trace is the Betti trace divided by $2\pi i$. Pullback multiplies the pairing by $d$, so
\[
 \left\langle\frac1m\pi^*\omega,\frac m d\pi^*\xi\right\rangle=1.
\]
The pullback image is two-dimensional, since the same nonzero pairing proves injectivity, and is the $f$-component. From the established complex decomposition
$\xi=-h\omega-c\bar\omega$, where $c=\pi/\operatorname{Area}(\Lambda)$, and $m\in\Q$, we get
\[
 \eta_f^{\rm alg}
 =-\frac{m^2}{d}h\omega_f
 -\frac{m^2}{d}c\overline{\omega_f}.
\]
This proves the exact scalar comparison.

Candelori \cite[Section 6, equation (6.5), Proposition 6.2, and equation (6.6), pp.27--30]{EllCandelori2014} represents the algebraic class by a second-kind differential $\phi$, regular away from the cusp at infinity, and constructs $F$ by
$dF=a\omega_f+b\overline{\omega_f}-\phi$. If
$\phi=\sum_n d(n)q^n\,dq/q$, then $d(n)\in\Q$, and equation (6.6) at weight two reads
\[
 [q^n]F^+=\frac{a\,b(n)-d(n)}{n}\qquad(n\ne0).
\]
All nonconstant coefficients thus belong to $\Q(a)$; conversely $b(1)=1$ gives
$a=[q]F^++d(1)$. With a rational constant term the generated field is exactly $\Q(a)=\Q(h)$. An algebraic constant term gives the stated algebraicity equivalence. This identifies the geometric construction through the actual pullback map, its cup-product normalization, and its coefficient formula.
\end{proof}

Candelori's final paragraph on p.30 explicitly identifies the weight-two converse with Katz's Question~4.0.8. The complete opened February~2025 version of Ehlen--Li--Schwagenscheidt \cite[Corollary 1.2 and Conjecture 1.5, pp.3--4]{EllEhlenLiSchwagenscheidt2024} proves the CM coefficient statement and states the converse shadow classification as a conjecture. At weight zero its algebraicity hypothesis is exactly the coefficient condition computed in A.4.23. Their Section~2.2 and Proposition~2.6 supply the explicit Hecke-character and binary-theta construction for the proved CM direction. No converse from that conjecture is used here, and the non-CM algebraicity question remains unresolved in this work.

\begin{result}{A.4.24}{The weight-two overconvergent criterion and its exact remaining equation}{Proved}
Let $f$ be the rational weight-two newform of an elliptic curve $E/\Q$, and let $p\ge5$ be a good ordinary prime. Write
\[
 T^2-b(p)T+p=(T-\beta)(T-\beta'),\qquad
 v_p(\beta)=0,\quad v_p(\beta')=1,
\]
and let $V(f)(q)=f(q^p)$ and $\theta=q\,d/dq$. Under the identification of the $f$-component with $H^1_{\dR}(E/\Q_p)$, the following three conditions are equivalent:
\[
 E\text{ has CM};\qquad
 f-\beta V(f)\in\operatorname{im}\theta;\qquad
 \Phi(\operatorname{Fil}^1)=\operatorname{Fil}^1.
\]
Here the image of $\theta$ is taken in the overconvergent spaces on the wide-open ordinary locus used in \cite[Theorem 3.3]{EllCandeloriCastella2017}, and $\Phi$ is crystalline Frobenius.

For a non-CM $E$, let $r$ be a weakly holomorphic weight-two form with zero residues whose class $\delta$ belongs to the $f$-component. Let $B(q)$ be its formal primitive, $\theta B=r$, and put $\varepsilon=\sum b(n)q^n/n$. There is a unique $\alpha\in\C_p$ for which
\[
 (B-\alpha\varepsilon)\bigm|\left(1-\frac{\beta'}pV\right)
\]
is overconvergent. It is characterized exactly by
\[
 (1-\Phi/\beta)(\delta-\alpha[f])=0.
\]
In particular, $\alpha=0$ precisely when $\delta$ belongs to the unit-root Frobenius line.
\end{result}
\begin{proof}
The cohomological comparison is explicit. Coleman's quotient theorem identifies the cokernel of $\theta$ with the overconvergent de Rham cohomology, and intertwines $U$ with crystalline Verschiebung. In weight two, crystalline Frobenius times Verschiebung is $p$ \cite[Theorems 5.1 and 5.4, pp.223 and 225]{EllColeman1996}. On cohomology $V=U^{-1}$, so
\[
 [Vf]=\Phi[f]/p.
\]
This also follows directly from pullback of the differential $f(q)\,dq/q$ under the canonical-subgroup Frobenius lift: its parameter map is $q\mapsto q^p$, and the differential factor is $p$. The restriction from the proper parabolic cohomology to the wide-open locus is injective \cite[equation (6), p.6]{EllCandeloriCastella2017}. Consequently
\[
 [f-\beta Vf]=0
 \quad\Longleftrightarrow\quad
 \Phi[f]=\frac p\beta[f]=\beta'[f].
\]
The modular pullback respects crystalline comparison, and $[f]$ generates the Hodge line. A Frobenius-stable integral Hodge line is divisible by $p$ under Frobenius; the determinant is $p$. Its eigenvalue must therefore be $\beta'$, proving the equivalence with the third condition. A.4.20 proves that this condition forces CM. For the reverse direction, Coleman's explicit CM theta construction \cite[Proposition 7.1, p.231]{EllColeman1996} produces an overconvergent primitive of $f-\beta Vf$. The rational weight-two converse is also the exact coefficient-field specialization of Emerton's proved Theorem~(1.3), pp.3--4 \cite{EllEmertonCM}; no higher-weight conjecture is needed here.

For the correction assertion, the identity $\theta V=pV\theta$ gives
\[
 \theta\left((B-\alpha\varepsilon)\bigm|\left(1-\frac{\beta'}pV\right)\right)
 =(1-\beta' V)(r-\alpha f).
\]
Its cohomology class is
$(1-\Phi/\beta)(\delta-\alpha[f])$. It vanishes exactly when the displayed formal primitive is overconvergent: the quotient theorem gives an overconvergent primitive, and the two primitives differ by a constant. The operator $1-\Phi/\beta$ has kernel the unit-root line and one-dimensional image. Its value on $[f]$ is nonzero, since the integral Hodge line cannot have unit-root eigenvalue. Thus exactly one scalar $\alpha$ solves the equation. Its vanishing is equivalent to $\delta$ lying in the kernel, as asserted.
\end{proof}

For the particular algebraic de Rham class that a putative algebraic $h$ supplies in A.4.23, $\delta=[\theta F^+]$ is a nonzero multiple of the complex antiholomorphic line, transported to $\C_p$ through a specified embedding of its coefficient field. The exact additional condition for an uncorrected overconvergent form is that this transported line equal the crystalline unit-root line. The formulas above establish that condition as an equality of two lines; they do not establish the equality from the algebraicity of $h$. Candelori--Castella's general correction theorems \cite[Theorems 5.3 and 5.5, pp.11--12]{EllCandeloriCastella2017} are therefore applicable in the non-CM rational weight-two case with their nonvanishing hypothesis verified, while the original H1 remains open in this work.


\begin{result}{A.4.25}{The exact defect between the complex and ordinary differential operators}{Proved}
On a fine modular curve, work locally with an algebraic symplectic frame $(\omega,\xi)$ of the universal elliptic de Rham bundle, with $\omega$ spanning the Hodge line and $\langle\omega,\xi\rangle=1$. Write
\[
 \nabla\omega=\omega\otimes A+\xi\otimes B,
 \qquad \operatorname{KS}(\omega^{\otimes2})=B.
\]
Here $B$ is a local generator of the cotangent bundle. Projection along the complement $L_a=\langle\xi+a\omega\rangle$, followed by Kodaira--Spencer, defines the weight-raising operator
\[
 D_a(f\omega^{\otimes k})
 =\left(\frac{df}{B}+k\left(\frac AB-a\right)f\right)
   \omega^{\otimes(k+2)}.
\]
For any two complements,
\[
 (D_a-D_b)(f\omega^{\otimes k})
   =k(b-a)f\omega^{\otimes(k+2)}.
\]
At a good ordinary point, let $L_{u_p}$ be the unit-root complement. For a specified embedding into $\C_p$ of a number field containing an algebraic $h$, agreement of the two operators on one section of nonzero weight and nonzero value is exactly
$\iota_p(h)=u_p$.
\end{result}
\begin{proof}
The symplectic definition of Kodaira--Spencer gives
$\langle\omega,\nabla\omega\rangle=B$, fixing its sign. Projection along $L_a$ sends $\omega$ to itself and $\xi$ to $-a\omega$. The Leibniz rule on $\omega^{\otimes k}$ therefore gives
\[
 \operatorname{pr}_a\nabla(f\omega^{\otimes k})
 =\bigl(df+kf(A-aB)\bigr)\omega^{\otimes k}.
\]
Applying the inverse Kodaira--Spencer identification proves the first formula. Subtraction proves the second. At the stated point, $k$ and $f$ are nonzero in characteristic zero; equality of the values is consequently equivalent to equality of the two complement parameters. A nonzero local section of the Hodge line supplies such a test section, so no global nonvanishing assumption on a particular Eisenstein series is needed.
\end{proof}

Mori's corrected final article \cite[Sections 3.3--3.5, pp.56--57]{EllMori1994} uses exactly Gauss--Manin, Kodaira--Spencer, and projection to construct these operators. Its Theorems~1 and~10 retain a CM elliptic curve with ordinary good reduction. Section~3.5 obtains equality of the complex and $p$-adic operators from the CM endomorphism decomposition. A.4.25 computes the equality required to apply that comparison in the present question. It has not been proved for a putative non-CM algebraic $h$. Section~0.4, p.47, explicitly corrects a constant in the 1990 and 1991 announcements; none of those uncorrected numerical expressions is used here.

\begin{result}{A.4.26}{The reduced theta coefficient field and its torsion translations}{Proved}
Let $k=\Q(g_2,g_3)$ for the fixed short Weierstrass equation, let $\sigma(u)$ be normalized by $\sigma'(0)=1$, and let $h$ be the actual Hodge splitting coefficient. The normalized reduced theta function and the associated Kronecker theta function are
\[
 \theta_h(u)=e^{-hu^2/2}\sigma(u),\qquad
 \Theta_h(u,v)=e^{-huv}\frac{\sigma(u+v)}{\sigma(u)\sigma(v)}.
\]
Their Taylor or Laurent coefficient fields over $k$ are exactly $k(h)$. Under the principal polarization, $\Theta_h$ represents the canonical meromorphic section of
\[
 m^*\mathcal O_E([0])\otimes
 p_1^*\mathcal O_E([0])^{-1}\otimes
 p_2^*\mathcal O_E([0])^{-1}
 \quad\hbox{on }E\times E.
\]
For every nonzero torsion point represented by $u_0\in\Lambda\otimes\Q$, the reduced translation
\[
 T_{u_0}(u)=
 e^{-c\bar u_0(u+u_0/2)}\theta_h(u+u_0),
 \qquad c=\frac{\pi}{\operatorname{Area}(\Lambda)},
\]
has coefficient field over $\Qb$ exactly $\Qb(h)$.
\end{result}
\begin{proof}
The formal equations $\wp'^2=4\wp^3-g_2\wp-g_3$, $\zeta'=-\wp$, and $\sigma'/\sigma=\zeta$, with their normalized leading terms, give
\[
 \sigma(u)=u-\frac{g_2}{240}u^5-\frac{g_3}{840}u^7+\cdots\in k[[u]].
\]
All subsequent coefficients are obtained by rational recursion. The coefficient of $u^3$ in $\theta_h$ is $-h/2$, and all coefficients lie in $k(h)$. This proves its exact field. Moreover,
\[
 \Theta_h(u,v)=\frac1u+\frac1v-h(u+v)
       +\text{terms of total degree at least three}.
\]
The same reasoning proves the exact Laurent coefficient field.

For $\lambda\in\Lambda$, the sigma transformation law and
$\eta_\lambda=h\lambda+c\bar\lambda$ give
\[
 \theta_h(u+\lambda)
 =\epsilon(\lambda)
   e^{c\bar\lambda(u+\lambda/2)}\theta_h(u).
\]
Together with $\theta_h'(0)=1$, this is the normalized reduced theta function of Bannai--Kobayashi \cite[Example 1.9, pp.10--11]{EllBannaiKobayashi2010}. Their area variable is $\operatorname{Area}(\Lambda)/\pi$, so its inverse is exactly our $c$. Taking the quotient under the addition map $m$ gives the stated section and line bundle, with residue one along $u=0$ \cite[Example 1.10, pp.11--12]{EllBannaiKobayashi2010}.

Extend $\lambda\mapsto\eta_\lambda$ $\Q$-linearly to $\Lambda\otimes\Q$. Direct multiplication gives
\[
 T_{u_0}(u)=e^{-hu^2/2}\,S_{u_0}(u),\qquad
 S_{u_0}(u)=e^{-\eta_{u_0}(u+u_0/2)}\sigma(u+u_0).
\]
Theorem~2.12, pp.26--27, of the same source applies to every algebraic elliptic curve, including those without CM, and gives algebraic coefficients for $S_{u_0}$. Its constant coefficient is nonzero because the torsion point is nonzero. If $s_j=[u^j]S_{u_0}$ and $t_j=[u^j]T_{u_0}$, then
$t_0=s_0\ne0$ and $t_2=s_2-hs_0/2$. Thus all $t_j$ lie in $\Qb(h)$ and $h=2(s_2-t_2)/s_0$, proving equality of fields.
\end{proof}

Bannai--Kobayashi's Proposition~2.14, pp.28--29, proves ordinary integrality by choosing an actual CM multiplication whose two eigenvalues have different $p$-adic valuations. Proposition~2.1 and Corollary~2.2, pp.20--21, likewise retain CM in their origin-algebraicity proof. The torsion-translation theorem used in A.4.26 has the broader hypothesis explicitly stated there; its exact Gaussian factor shows what it transfers. The theorem has not supplied ordinary bounded denominators from $h\in\Qb$, and H1 remains unchanged.


\begin{result}{A.4.27}{The exact G-function input for the elliptic primitive}{Proved}
Let $K$ be a number field, let $E:y^2=P(x)=4x^3-g_2x-g_3$ be smooth over $K$, and let $Q=(x_0,y_0)\in E(K)$ satisfy $y_0\ne0$. For every $a\in K$, the germ
\[
 F_a(t)=-\int_0^t\frac{x_0+s+a}{\sqrt{P(x_0+s)}}\,ds,
 \qquad \sqrt{P(x_0)}=y_0,
\]
is a $G$-function over $K$. In the coordinate $x=x_0+t$ it satisfies
\[
 2P(x)(x+a)F_a''+
 \bigl(P'(x)(x+a)-2P(x)\bigr)F_a'=0.
\]
For some integer $C\ge1$, the first $n$ Taylor coefficients have a common denominator dividing
\[
 (4C)^{n+1}\operatorname{lcm}(1,\ldots,n).
\]
If $x=\wp(u)$ and $x_0=\wp(u_0)$, the exact analytic expression is
\[
 F_a(x-x_0)=\zeta(u)-au-\zeta(u_0)+au_0.
\]
At $a=h$, its additive monodromy along an elliptic loop of period $\lambda$ is $c\bar\lambda$, and is nonzero whenever $\lambda\ne0$.
\end{result}
\begin{proof}
Write $P(x_0+t)=y_0^2(1+b_1t+b_2t^2+b_3t^3)$. The expansion
\[
 (1+z)^{-1/2}=
 \sum_{m\ge0}\frac{(-1)^m\binom{2m}{m}}{4^m}z^m
\]
computes the derivative of $F_a$. Choose $C$ clearing the denominators in $\mathcal O_K$ of the $b_i$, $(x_0+a)/y_0$, and $1/y_0$. The coefficient $e_n=[t^n]F_a'$ then satisfies $(4C)^{n+1}e_n\in\mathcal O_K$: each summand uses at most $n$ of the $b_i$ and one further cleared factor. Integration divides $e_{j-1}$ by $j$, proving the denominator bound. The least common multiple grows at most exponentially, and each complex conjugate of the algebraic derivative is analytic on a nonzero disc at zero. Cauchy's estimate therefore bounds all conjugates of the coefficients exponentially. Together with the displayed linear differential equation, these are precisely the $G$-function conditions. That equation follows by differentiating $F_a'=-(x+a)/\sqrt{P(x)}$ and clearing denominators.

The uniformization identity follows from $d\zeta=-x\,du$ and $du=dx/y$, with the constant fixed at $u_0$. Translation by $\lambda$ changes it by $\eta_\lambda-a\lambda$. The established identity $\eta_\lambda=h\lambda+c\bar\lambda$ proves the last assertion.
\end{proof}

The current arithmetic-holonomy theorem of Calegari--Dimitrov--Tang \cite[Theorem 2.5.1, pp.16--17]{EllCDT2024} imposes an explicit denominator template and a strict inequality comparing a holomorphic uniformizing map's conformal radius with its denominator exponent. Its numerator is an explicit Bost--Charles boundary double integral. A.4.27 verifies the elementary $G$-function input for every algebraic correction, while also computing the actual monodromy at $h$. A quantitative analytic improvement or a further product-formula relation specific to $h$ has not been established here. The source does not infer ordinary bounded denominators from the $G$-function condition. The related 2026 ICM article \cite{EllCDT2026} was also opened; no unverified application of its effective bounds is used.

\begin{result}{A.4.28}{Horizontal Hodge jets and an explicit CM denominator obstruction}{Proved}
In the Legendre family
\[
 E_t:y^2=4x(x-1)(x-t),\qquad
 \omega=dx/y,\quad \eta=x\,dx/y,
\]
the Gauss--Manin connection matrix, with columns giving the derivatives of the basis, is
\[
 C(t)=\begin{pmatrix}A&A\\ B&-A\end{pmatrix},\qquad
 A=\frac1{2(1-t)},\qquad B=-\frac1{2t(1-t)}.
\]
At $t_0=1/2$, the antiholomorphic class is $\eta-\omega/2$, and the short Weierstrass equation has $g_2=1$, $g_3=0$, and $h=0$. Its horizontal extension $F(s)\omega+G(s)\eta$, with $s=t-1/2$, is
\[
 G(s)=\sum_{m\ge0}
 \frac{\prod_{j=0}^{m-1}(4j+1)^2}{(2m)!}s^{2m},
 \qquad
 F(s)=\left(\frac12-2s^2\right)G'(s)
       -\left(\frac12+s\right)G(s).
\]
For every prime $p\equiv3\pmod4$, the coefficient $[s^{p+1}]G$ has $p$-adic valuation $-1$. Thus this actual algebraic CM Hodge splitting has horizontal Taylor coefficients with infinitely many denominator primes, and fails the $\omega(p)$-integrality condition of Lam--Litt.

More generally, at every algebraic nonsingular Legendre parameter and every nonzero algebraic initial vector in a positive symmetric power of its de Rham cohomology, the horizontal extension fails both integrality conditions of \cite[Theorem 4.1.1, p.21]{EllLamLitt2025}.
\end{result}
\begin{proof}
Differentiation at fixed $x$ gives
$\partial_t\omega=\omega/(2(x-t))$ and
$\partial_t\eta=\omega/2+t\partial_t\omega$.
For $R=x(x-1)/y$, direct differentiation gives
\[
 d_xR=\left(\frac{x-t}{2}
       -\frac{t(t-1)}{2(x-t)}\right)\omega.
\]
Modulo exact forms, these identities give the stated matrix, including its signs.
At $t=1/2$, the translation $x=X+1/2$ gives $y^2=4X^3-X$; hence $\eta-\omega/2=X\,dX/y$ is the antiholomorphic class by A.4.4.

Horizontal coefficient vectors satisfy $(F,G)'=-C(F,G)$. Eliminating $F$ gives
\[
 t(1-t)G''+(1-2t)G'-G/4=0,
 \qquad G(0)=1,\quad G'(0)=0
\]
where the last two values are at $s=0$. Equivalently,
$(1-4s^2)G''-8sG'-G=0$. Its coefficient recurrence is
\[
 [s^{2m+2}]G=\frac{(4m+1)^2}{(2m+2)(2m+1)}[s^{2m}]G,
\]
proving the formula and the expression for $F$.

For $m=(p+1)/2$, every numerator factor $4j+1$ is positive and below $2p$. None equals $p$, since $p\equiv3\pmod4$. Thus the numerator is a $p$-adic unit, whereas $(2m)!=(p+1)!$ has valuation one. There are infinitely many primes congruent to three modulo four: if they formed a finite list, a prime divisor congruent to three modulo four of four times their product minus one would be new. Consequently infinitely many denominator primes occur. If $\omega(p)/p\to\infty$, then $p+1\le\omega(p)$ for every sufficiently large prime; the exhibited coefficients also exclude $\omega(p)$-integrality.

For the general statement, take $R=\mathcal O_K[1/N]$ large enough to contain the algebraic initial data and to make $t_0$ and $1-t_0$ units. The Legendre family over $R[t,1/(t(1-t))]$ is smooth, projective, and nonisotrivial. Its $j$-function is nonconstant. The nonzero initial vector and every positive symmetric power therefore satisfy exactly the hypotheses of Lam--Litt's Theorem~4.1.1. This gives the asserted specialization of that proved theorem.
\end{proof}

The relationship to the complex splitting is explicit. In a locally constant Betti trivialization, the antiholomorphic period vector depends antiholomorphically on the base coordinate. Freezing its conjugate coordinate at the starting point gives its horizontal holomorphic continuation. Thus an algebraic $\xi+h\omega$ supplies the algebraic initial vector in A.4.28, with the scalar normalization fixed at that point. It does not supply integral horizontal jets; the displayed CM example disproves that inference. Lam--Litt's theorem consequently applies to this continuation with its stated hypotheses, while the original algebraicity question for $h$ remains unproved.

\begin{result}{A.4.29}{Division-polynomial data preserve the Gaussian parameter ambiguity}{Proved}
For any $a\in\C$, put $\theta_a(u)=e^{-au^2/2}\sigma(u)$ and
$\eta^{(a)}_\lambda=\eta_\lambda-a\lambda$. The normalized division polynomial $F_n$ satisfies
\[
 F_n(\wp(u),\wp'(u))
 =\frac{\sigma(nu)}{\sigma(u)^{n^2}}
 =\frac{\theta_a(nu)}{\theta_a(u)^{n^2}}.
\]
For every $u_0\in\Lambda\otimes\Q$, the normalized torsion constant likewise satisfies
\[
 e^{-\eta^{(a)}_{u_0}u_0/2}\theta_a(u_0)
 =e^{-\eta_{u_0}u_0/2}\sigma(u_0).
\]
Thus the map from the parameter $a$ to all division-polynomial quotients and these torsion constants is constant. The same quotient identities hold as formal series over every characteristic-zero coefficient field, including a $p$-adic field.
\end{result}
\begin{proof}
The Gaussian factors in the first quotient cancel because $(nu)^2=n^2u^2$. The sigma transformation law shows that the remaining quotient is elliptic. Its divisor is $[n]^*[O]-n^2[O]$, and its leading term is $n u^{1-n^2}$. As $t=-2x/y=u+O(u^2)$, this is exactly the normalization
\[
 \left.\frac{t^{n^2}F_n}{[n]^*t}\right|_O=1
\]
used in \cite[p.2]{EllKatagiriKobayashi2026}. Divisor and leading term uniquely determine the rational function, proving the first identity. Substitution of $\eta^{(a)}_{u_0}=\eta_{u_0}-au_0$ proves the second. The Gaussian cancellation is an identity of formal exponential series, proving its characteristic-zero formal version.
\end{proof}

Katagiri--Kobayashi's September~2026 primary preprint \cite[Theorem 1.3 and Corollary 1.4, pp.3--4]{EllKatagiriKobayashi2026} determines certain $p$-power subsequential limits of division-polynomial values in both ordinary and supersingular good reduction. Its algebraic sigma value is exactly the normalized torsion constant above \cite[Definition 2.3, Proposition 2.5, and Theorem 2.7, pp.6--7]{EllKatagiriKobayashi2026}. A.4.29 identifies the actual noninjective map through which these data pass. It does not identify the Gaussian parameter with the canonical complex value $h$, nor does it exclude a different construction using further data. The remaining H1 gate is unchanged.

\begin{result}{A.4.30}{The Hodge parameter uniquely minimizes the quadratic growth type}{Proved}
Let $c=\pi/\operatorname{area}(\Lambda)>0$ and, for $a\in\C$, put
\[
 \theta_a(u)=e^{-au^2/2}\sigma(u),\qquad
 M_a(r)=\max_{|u|\le r}|\theta_a(u)|.
\]
Then
\[
 \log M_a(r)=\frac{c+|h-a|}{2}\,r^2+O(r)
 \qquad(r\longrightarrow+\infty).
\]
In particular, every $\theta_a$ has order two, its quadratic type is
$(c+|h-a|)/2$, and $a=h$ is the unique parameter of minimal type.
This gives an exact analytic property specific to the actual Hodge
parameter, in addition to the coefficient-field identity in A.4.26.
\end{result}
\begin{proof}
The reduced transformation law proved in A.4.26 gives
\[
 |\theta_h(u+\lambda)|
 =e^{c\operatorname{Re}(\bar\lambda u)+c|\lambda|^2/2}
   |\theta_h(u)|.
\]
Consequently
\[
 B(u):=|\theta_h(u)|e^{-c|u|^2/2}
\]
is a continuous $\Lambda$-periodic function. It is bounded above,
and $B(u_0)>0$ for any fixed $u_0\notin\Lambda$. Since
\[
 |\theta_a(u)|=
 B(u)\exp\!\left(\frac{c|u|^2+
                 \operatorname{Re}((h-a)u^2)}2\right),
\]
the inequality
$\operatorname{Re}((h-a)u^2)\le |h-a||u|^2$
proves the required upper bound, with an $O(1)$ error.

Choose an angle $\phi$ with
$\operatorname{Re}((h-a)e^{2i\phi})=|h-a|$; when $a=h$ any
angle will do. The translated lattice $u_0+\Lambda$ has a finite
covering radius $R$. For every sufficiently large $r$, choose
$u_r\in u_0+\Lambda$ at distance at most $R$ from
$(r-2R)e^{i\phi}$. Then $|u_r|\le r$, $B(u_r)=B(u_0)>0$, and
$u_r=re^{i\phi}+O(1)$. The real quadratic polynomial
$c|u|^2+\operatorname{Re}((h-a)u^2)$ therefore has value
$(c+|h-a|)r^2+O(r)$ at $u_r$. Evaluation there proves the lower
bound. The assertions about order, type, and the unique minimizer
follow because $c>0$.
\end{proof}

Two opened arithmetic growth results were checked against this exact
calculation. Herblot's \cite[Theorem 6.1, pp.31--32]{EllHerblot2012}
uses the order of a connected meromorphic uniformization, together
with arithmeticity parameters and residue degrees; its displayed
inequality does not use the quadratic type. The scalar theta
functions in A.4.30 all have order two, so their minimal-type
property alone supplies no improvement to an order estimate using
these functions.
Fonseca's \cite[Theorem 1.2, p.648]{EllFonseca2019} instead requires
a formal morphism over $\mathcal O_K[[q]]$, analytic continuations at
every embedding with product of radii one, moderate growth, and
ZL-density. Its conclusion concerns nonzero arguments.
For the actual horizontal germ in A.4.28, its center has argument
zero and the explicitly exhibited denominator primes prevent an
integral formal morphism in that frame over every
$\mathcal O_K[1/N]$ with fixed $N$. These statements check the
specified map and its hypotheses; they do not exclude another
map with independently verified arithmetic estimates.
No product-formula estimate combining the minimal type in A.4.30
with the required finite-place bounds has been established here.
The non-CM algebraicity gate for $h$ remains unresolved.

\begin{result}{A.4.31}{The corrected elliptic vector field at every good prime}{Proved}
Let $k$ be a perfect field of characteristic $p>3$, and let
\[
 E/k:\quad y^2=P(x)=4x^3-g_2x-g_3
\]
be smooth. Define the two explicit coefficients
\[
 H_p=[x^{p-1}]P(x)^{(p-1)/2},\qquad
 K_p=[x^{p-2}]P(x)^{(p-1)/2}.
\]
In the rational universal-extension chart of A.4.14, for $a\in k$ put
\[
 D_a=y\partial_x+
       \left(6x^2-\frac{g_2}{2}\right)\partial_y
       -(x+a)\partial_t,\qquad V=-\partial_t.
\]
The restricted $p$-th power, defined by the $p$-fold iterate of a
derivation, is
\[
 \boxed{\quad D_a^p=H_pD_a-(K_p+aH_p)V.\quad}
\]
Consequently the line $kD_a$ is stable under restricted $p$-th power
exactly when $K_p+aH_p=0$. In ordinary reduction this determines
the unique value $a=-K_p/H_p$. In supersingular reduction no
line $kD_a$ is stable under restricted $p$-th power.
\end{result}
\begin{proof}
Use the classical $p^{-1}$-linear Cartier operator on rational
differentials over the perfect field $k$. Its coefficient convention is
\[
 C(f^px^m\,dx)=0\quad(m\not\equiv p-1\pmod p),\qquad
 C(f^px^{p-1}\,dx)=f\,dx.
\]
Since $y^{p-1}=P(x)^{(p-1)/2}$, the forms of A.4.1 satisfy
\[
 \omega=\frac{P(x)^{(p-1)/2}\,dx}{y^p},\qquad
 \xi=\frac{xP(x)^{(p-1)/2}\,dx}{y^p}.
\]
Both numerator degrees are less than $2p-1$. There is therefore
exactly one possible extracted exponent, giving
\[
 C(\omega)=H_p^{1/p}\omega,\qquad
 C(\xi)=K_p^{1/p}\omega.
\]
These roots and the displayed semilinearity retain the Frobenius
convention explicitly.

For a closed rational one-form $\alpha$ and a derivation $D$,
Cartier's contraction identity is
\[
 \alpha(D^p)-D^{p-1}\bigl(\alpha(D)\bigr)
       =\bigl(C(\alpha)(D)\bigr)^p.
\]
This is \cite[Proposition 1.2.4, equation (1.2.4.5), p.24]{EllOgus1994}
after writing its relative-Frobenius pairing in the preceding
classical semilinear convention; its inverse-Cartier verification
is on p.25.
The closed forms
\[
 \omega,\qquad \beta=-dt-\xi
\]
are a rational coframe on $k(E)(t)$, and
\[
 \omega(D_a)=1,\quad \beta(D_a)=a,\qquad
 \omega(V)=0,\quad \beta(V)=1.
\]
Cartier kills $dt$, hence $C(\beta)=-K_p^{1/p}\omega$.
Applying the contraction identity gives
\[
 \omega(D_a^p)=H_p,\qquad \beta(D_a^p)=-K_p.
\]
The coframe determines a derivation uniquely, proving the boxed
identity with its stated sign. Equivalently, if
$\delta=y\partial_x+(6x^2-g_2/2)\partial_y$, the same calculation
on $E$ gives
\[
 \delta^p=H_p\delta,\qquad
 \delta^{p-1}x=H_px-K_p.
\]
Its application to $t$ gives
$D_a^pt=-H_px+K_p$, which checks the vertical component directly.

For completeness, $H_p$ and $K_p$ cannot both vanish on a smooth
elliptic curve. The conjugate-filtration sequence is
\[
 0\longrightarrow H^1(E^{(p)},\mathcal O_{E^{(p)}})
 \longrightarrow H^1_{\mathrm{dR}}(E/k)
 \longrightarrow H^0(E^{(p)},\Omega^1_{E^{(p)}/k})
 \longrightarrow0.
\]
The final map is induced by Cartier, with the twist interpreted
as above. The meromorphic form $\xi$ gives an actual de Rham
hypercohomology class: for the parameter $z=-2x/y$ at $O$,
\[
 \omega=(1+O(z^4))\,dz,\qquad
 \xi=(z^{-2}+O(z^2))\,dz.
\]
Its principal part is $d(-1/z)$. Removing this exact principal
part on a neighborhood of $O$ gives regular local representatives
and the overlap function $1/z$. Cartier kills their exact
transition. The residue computation
$\operatorname{Res}_O((1/z)\omega)=1$ gives the cup pairing
$\langle\omega,\xi\rangle=1$ in the convention of A.4.1.
Thus these two classes form a basis. Surjectivity onto the
one-dimensional final term shows that their Cartier images
cannot both vanish. Ordinary reduction is characterized by
$H_p\ne0$; in the supersingular case $H_p=0$ the nonvanishing
therefore gives $K_p\ne0$. This proves all the assertions about
the lines $kD_a$.
\end{proof}

At $p=5$ the coefficients are $H_5=2g_2$ and $K_5=2g_3$.
The formula consequently reads
\[
 D_a^5=2g_2D_a+(3g_3-2g_2a)V,
\]
agreeing with the earlier exact calculation.
The script \texttt{check\_cartier\_all\_primes.py} independently
checks the polynomial identities for generic $g_2,g_3$ at
$p=5,7,11,13,17,19$; the argument above proves the all-prime
statement. For an algebraic proposed value of the complex
coefficient $h$, the ordinary finite-place equation is now
explicitly $K_p+\overline h\,H_p=0$, where $\overline h$ denotes
reduction at the specified good prime and requires $h$ to be
integral there. No implication from the algebraicity of the
actual complex $h$ to these equations has been established.
The supersingular CM calculation of A.4.14 continues to exclude
an unrestricted assertion of $p$-closure. The main H1 gate is
unchanged.

\begin{result}{A.4.32}{A density bound for the explicit restricted-Lie congruences}{Proved}
Let $K$ contain $g_2,g_3,a$, where $a$ is algebraic and
$E:y^2=4x^3-g_2x-g_3$ is smooth, and put $d=[K:\Q]$.
Omit finitely many primes so that all data have good integral
reduction and their residue characteristics exceed three.
For a prime $v$ of $K$ above the rational prime $p$, put
\[
 w_v=[K_v:\Q_p]\frac{\log p}{p-1},\qquad
 \beta_K(a)=\liminf_{R\to\infty}
 \frac{\displaystyle\sum_{\substack{v\mid p,\ p\le R\\
                         K_p+aH_p=0\ {\rm in}\ k(v)}}w_v}
      {\displaystyle\sum_{v\mid p,\ p\le R}w_v}.
\]
Here $H_p,K_p$ are the reductions of the coefficients in A.4.31.
Then
\[
 \boxed{\qquad \beta_K(a)\le 1-\frac1{4d}.\qquad}
\]
In particular, the bound over $\Q$ is $3/4$. Equivalently,
the complementary weighted upper density is at least $1/(4d)$.

The analytic input is uniform in the correction parameter:
the curve $\Theta_a(u)=\exp(u,au)$ has an order-two
uniformization in an explicit projective compactification.
For the ample bundle constructed below, the least quadratic
type of a nonzero holomorphic growth section is $Nc/2$ for
every $a$, where $c=\pi/\operatorname{area}(\Lambda)$.
\end{result}
\begin{proof}
The additive torsor $G=E^\dagger\to E$ determines an algebraic
extension
\[
 0\longrightarrow\mathcal O_E\longrightarrow\mathcal V
 \longrightarrow\mathcal O_E\longrightarrow0.
\]
Explicitly, an additive torsor cocycle $b_{ij}$ gives the transition
matrices $\left(\begin{smallmatrix}1&b_{ij}\\0&1\end{smallmatrix}\right)$.
In the projective bundle $X=\mathbb P_E(\mathcal V)$ of lines,
the complement of $\mathbb P_E(\mathcal O_E)$ is $G$:
the affine fiber coordinate is $s=v_0/v_1$ on $v_1\ne0$.
This constructs a smooth projective compactification over $K$
containing the identity.

At a fixed embedding $K\hookrightarrow\C$, A.4.14 gives the
analytic cocycle
\[
 (u,v_0,v_1)\sim
 (u+\lambda,v_0+\eta_\lambda v_1,v_1).
\]
Thus the required map into the compactification is exactly
\[
 \Theta_a:\C\longrightarrow G(\C)\subset X(\C),\qquad
 u\longmapsto[u;(au,1)].
\]
It uniformizes the formal leaf of $L_a$ at the identity.
Put $\mathcal B=\mathcal O_E([O])$, and choose $N$ large enough
that
\[
 \mathcal L=\mathcal O_{\mathbb P(\mathcal V)}(1)
                  \otimes\pi^*\mathcal B^{\otimes N}
\]
is ample. Give $\mathcal V$ a smooth Hermitian metric. The
tautological vector $v_a(u)=(au,1)$ is a nowhere-zero holomorphic
section of $\Theta_a^*\mathcal O_{\mathbb P(\mathcal V)}(-1)$.
Writing $u=w+\lambda$ with $w$ in a compact fundamental
parallelogram identifies this vector with
$(au-\eta_\lambda,1)$ over $w$. Since
$|\eta_\lambda|\le C|\lambda|$, compactness and positivity of
the metric give
\[
 C_1\le\|v_a(u)\|\le C_a(1+|u|).
\]
For the lower bound, the fixed quotient
$\mathcal V\to\mathcal O_E$ sends this vector to $1$ and has
bounded norm on the compact base.

The reduced theta multiplier in A.4.26 presents
$\mathcal B$ with a holomorphic frame $e_h$ on its pullback
to $\C$. Give this frame norm
$\|e_h(u)\|=e^{-c|u|^2/2}$.
The multiplier has modulus
$e^{c\operatorname{Re}(\bar\lambda u)+c|\lambda|^2/2}$;
the ratio of these frame norms is its inverse, so the metric
descends to $\mathcal B$.
The nowhere-zero holomorphic section
\[
 s_a(u)=v_a(u)\otimes e_h(u)^{-N}
       \quad\hbox{of }\Theta_a^*\mathcal L^{-1}
\]
therefore satisfies
\[
 C_1e^{Nc|u|^2/2}\le\|s_a(u)\|
      \le C_a(1+|u|)e^{Nc|u|^2/2}.                 \tag{*}
\]
This verifies precisely the growth-section definition
\cite[Definition 5.1, p.29; Definitions 5.4--5.5, p.31]{EllHerblot2012}
with $M=\mathbb P^1$, puncture $\infty$, and order $\rho=2$.
The metric uses the actual complex $h$ of $E$, whose algebraicity
is not needed for its construction.

Every other nonzero holomorphic section of this pulled-back
inverse line bundle is $f(u)s_a(u)$ for a nonzero entire
function $f$. If $f$ vanishes to order $m$ at zero, Jensen's
formula for $f(u)/u^m$ and the lower bound in $(*)$ give
\[
 \log\max_{|u|=r}\|f(u)s_a(u)\|
       \ge Nc\,r^2/2+m\log r+O(1).
\]
The upper bound in $(*)$ attains quadratic type $Nc/2$.
This proves the stated optimal growth constant for every $a$.

The Zariski closure of $\Theta_a(\C)$ in $G$ is a connected
algebraic subgroup: multiplication and inversion preserve
the closure of this connected analytic subgroup.
Its projection to $E$ is surjective. A one-dimensional closure
would have Lie algebra $L_a$, contrary to A.4.15. Hence the
closure has dimension two. The identity theorem transfers
vanishing on the formal leaf to vanishing on the analytic curve,
so the formal leaf has the same closure.

The invariant line is a nonsingular algebraic foliation on
the open $G\subset X$. By A.4.31 its reduction is closed under
$p$-th powers exactly at the primes used in the numerator
defining $\beta_K(a)$. Herblot's density and arithmeticity
result \cite[Definition 3.5 and Proposition 3.6, pp.20--21]{EllHerblot2012}
makes the formal leaf $(1-\beta_K(a))$-arithmetic.

Regard $X$ as an integral projective $\Q$-variety through
$\operatorname{Spec}K\to\operatorname{Spec}\Q$. It is smooth,
and still has dimension two, since this base map is finite
and \'etale. Its identity is now a closed point with residue
field $K$ and degree $d$.
Arithmeticity survives passage to a Galois closure by
\cite[Lemma 2.6]{EllHerblot2012}; the conjugacy argument
in \cite[Definition 4.6 and its remark, p.28]{EllHerblot2012}
then gives the required arithmeticity at every residue-field
embedding. The analytic definition requires one geometric
point above each closed point, so the chosen embedding
$K\hookrightarrow\C$ supplies the required uniformization.
A $\Q$-closed subset containing the dense chosen complex
component contains all its Galois conjugates. Thus the
$\Q$-Zariski closure has dimension $r=2$, even though
$X_\C$ has components indexed by the embeddings of $K$.

Herblot's \cite[Theorem 6.1, pp.31--32]{EllHerblot2012}
now applies with one closed point of degree $d$, order
$\rho=2$, and arithmeticity parameter $1-\beta_K(a)$.
Its inequality is
\[
 \frac1{(1-\beta_K(a))d}\le
       \frac{2}{2-1}\,2=4.
\]
When necessary, use a slightly larger positive arithmeticity
parameter and let it decrease to $1-\beta_K(a)$.
This proves the density bound, including exclusion of
$\beta_K(a)=1$.
\end{proof}

The density in A.4.32 is the specified weighted lower density.
The complementary conclusion is correspondingly about an upper
density. It strengthens A.4.15's infinite-prime failure for the
actual algebraic line, with the residue-field degree retained.
The scalar-theta minimization of A.4.30 supplies no smaller growth
constant in this particular compactified uniformization, since
$(*)$ and Jensen's formula compute the same minimum for every
correction parameter. No density of the congruences
$K_p+hH_p=0$ has been deduced from the algebraicity of the actual
complex $h$. H1 remains unresolved.

\begin{result}{A.4.33}{The conjugate elliptic quotient and its differential scalar}{Proved}
Let $E/\Qb$, $h$, and $c=\pi/\operatorname{area}(\Lambda)>0$
have the preceding normalizations. The algebraic conjugate curve
$\overline E$ has differential-normalized lattice $\bar\Lambda$.
There is a surjective analytic group homomorphism
\[
 \Phi_h:E^\dagger(\C)\longrightarrow\overline E(\C),
 \qquad [(u,s)]\longmapsto[(s-hu)/c]\pmod{\bar\Lambda}.
\]
Its kernel is $\exp(L_h)$, and $u\mapsto[(u,hu)]$ is an
injective uniformization of that kernel. On the maximal compact
subgroup, it satisfies
$\Phi_h(\kappa([u]))=[\bar u]$.
The homomorphism $\Phi_h$ is never algebraic over $\C$.

Writing $v$ for the holomorphic uniformizing coordinate of
$\overline E$, its differential is
\[
 \Phi_h^*(dv)=c^{-1}(ds-h\,du),
 \qquad d\Phi_h=(-h/c,\ 1/c)
\]
in the algebraic coordinates of A.4.14. The pair $(h,c)$
cannot belong to $\Qb^2$. In particular, whenever $h$ is
algebraic, $c$ is transcendental.

For the actual CM example $(g_2,g_3,h)=(30,28,1/2)$,
$\Phi_h$ is consequently an explicit nonalgebraic analytic
homomorphism between algebraic groups whose differential has
algebraic projective class $[-1/2:1]$.
\end{result}
\begin{proof}
The identity
$\eta_\lambda-h\lambda=c\bar\lambda$
shows that replacing $(u,s)$ by $(u+\lambda,s+\eta_\lambda)$
changes $(s-hu)/c$ by $\bar\lambda$. This proves descent.
The choice $(u,s)=(0,cz)$ proves surjectivity.
If the image of $[(u,s)]$ is zero, choose $\lambda\in\Lambda$
with $s-hu=c\bar\lambda$. Subtracting
$(\lambda,\eta_\lambda)$ gives
\[
 s-\eta_\lambda-h(u-\lambda)
 =(s-hu)-(\eta_\lambda-h\lambda)=0.
\]
The resulting representative belongs to $L_h$, proving the
kernel equality. Conversely, $L_h$ maps to zero.
If $(u,hu)$ belongs to $\Lambda^\dagger$, then $u=\lambda$
and $c\bar\lambda=\eta_\lambda-h\lambda=0$, hence $\lambda=0$.
This proves injectivity. Substitution of
$\kappa([u])=[(u,hu+c\bar u)]$ proves the compact-subgroup
identity.

On the algebraic vertical subgroup $\Ga$, the restriction is
$s\mapsto[s/c]\in\C/\bar\Lambda$ and has derivative $1/c\ne0$.
An algebraic map $\mathbb A^1_\C\to\overline E$ extends at
infinity to $\mathbb P^1_\C$ by properness. A nonconstant map
from $\mathbb P^1$ to an elliptic curve would contradict
Riemann--Hurwitz, which would give
$-2=0+\text{nonnegative ramification}$.
Thus the restriction, and therefore $\Phi_h$, is not algebraic.

The torsion points of $E^\dagger$ are Zariski dense.
Indeed, $E^\dagger[n]\to E[n]$ is an isomorphism because
multiplication by $n$ is invertible on the additive kernel.
The algebraic closure of all torsion is a subgroup $H$
projecting onto $E$. Its identity component also projects
onto $E$: otherwise finitely many components would have
finite image, although $E[n]$ has unbounded order.
If $\dim H=1$, this would be the forbidden one-dimensional
algebraic subgroup in A.4.15. Hence $H=E^\dagger$.

Huber--W\"ustholz's dense-torsion analytic-homomorphism
lemma \cite[Lemma 6.5, pp.45--46 of the opened manuscript]{EllHuberWustholz2022}
states that an analytic homomorphism between connected
commutative algebraic groups over $\Qb$ is algebraic when
its differential is algebraic and source torsion is
Zariski dense. Both groups here are defined over $\Qb$,
and the differential is the displayed row.
If $h,c$ were both algebraic, the lemma would make $\Phi_h$
algebraic, contrary to the preceding proof.
Thus $(h,c)\notin\Qb^2$. The CM value $h=1/2$ was proved
in A.4.4, and its stated projective differential follows
by multiplying the row by $c$.
\end{proof}

The exact change of differential normalization can also be
written down. Put
\[
 E_c^{\mathrm{an}}=\C/(c\bar\Lambda),\qquad
 \widetilde\Phi_h([(u,s)])=[s-hu].
\]
This presentation has differential $(-h,1)$, and its
short Weierstrass coefficients in the differential $dz$ are
\[
 g_2(E_c)=c^{-4}\bar g_2,\qquad
 g_3(E_c)=c^{-6}\bar g_3.
\]
The actual analytic isomorphism and its differential are
\[
 \mu_c:E_c^{\mathrm{an}}\longrightarrow\overline E^{\mathrm{an}},
 \quad[z]\longmapsto[z/c],\qquad
 \mu_c^*(dv)=dz/c,\qquad
 \Phi_h=\mu_c\circ\widetilde\Phi_h.
\]
In Weierstrass coordinates it is $(X,Y)\mapsto(c^2X,c^3Y)$,
by the scaling identity
$\wp_{c\bar\Lambda}(z)=c^{-2}\wp_{\bar\Lambda}(z/c)$.
When $h$ is algebraic, A.4.33 makes $c$ transcendental.
Every nonzero coefficient among $c^{-4}\bar g_2$ and
$c^{-6}\bar g_3$ is then transcendental; algebraicity of
one would make a positive integral power of $c$, and hence
$c$ itself, algebraic. At least one is nonzero.
Although $j(E_c)=\overline{j(E)}$ is algebraic, passage to
the algebraic model $\overline E$ uses precisely $\mu_c$
and restores the factor $1/c$ in the differential.

The consulted manuscript is the complete primary version~8
of 21~April~2022. Its Remark~6.4, p.45, explicitly corrects
an earlier statement which omitted a vector-group factor;
the exponential map supplies the stated exception.
The application above uses the dense-torsion Lemma~6.5
and retains the actual differential hypothesis.
Algebraicity of the projective row alone cannot replace it,
as the explicit CM example demonstrates.
The resulting transcendence of $c$ when $h$ is algebraic
does not distinguish CM from a putative non-CM algebraic
splitting. The main H1 gate remains unchanged.

\subsection{Second-kind period lattices and the CM quotient}

\begin{result}{A.4.34}{The exact quotient lattice and its CM case}{Proved}
Let $E/\Qb$, $\omega_i,\eta_i,h,c,A$ have the conventions of A.4.1,
and put
\[
 \mu_i(a)=\eta_i-a\omega_i,\qquad
 \Lambda_a=\Z\mu_1(a)+\Z\mu_2(a).
\]
For $a\in\C$, this is a full lattice in $\C$ exactly when
$|h-a|\ne c$. Off that circle its basis ratio satisfies
\[
 r_a=\frac{\eta_2-a\omega_2}{\eta_1-a\omega_1},
 \qquad
 \operatorname{Im}r_a
 =\frac{A\bigl(|h-a|^2-c^2\bigr)}
        {|\eta_1-a\omega_1|^2}.
\]
There is a holomorphic map
\[
 F_a:E(\C)\setminus\{O\}\longrightarrow \C/\Lambda_a,
 \qquad [u]\longmapsto[\zeta(u)-au],
 \qquad F_a^*(dz)=-(x+a)\omega.
\]
It has no holomorphic extension at $O$ and is not algebraic.
For algebraic $a$ and a full lattice $\Lambda_a$, the complex elliptic
curve $\C/\Lambda_a$ has CM if and only if $E$ has CM and $a=h$.
\end{result}

\begin{proof}
The period formula of A.4.1 gives the exact real-linear map
\[
 T_a:\C\longrightarrow\C,\qquad
 z\longmapsto(h-a)z+c\bar z,
 \qquad T_a(\Lambda)=\Lambda_a.
\]
Writing $h-a=b_1+ib_2$ in real coordinates gives its matrix
\[
 \begin{pmatrix}b_1+c&-b_2\\ b_2&b_1-c\end{pmatrix},
 \qquad \det_{\R}T_a=|h-a|^2-c^2.
\]
Off the stated circle this is an invertible real-linear map, so it
carries the full lattice $\Lambda$ to a full lattice. On the circle it
has rank one, because $c>0$, and its image cannot be a full lattice in
$\C$. No claim of discreteness is made for that rank-one-line image.
Off the circle $\mu_1(a)\ne0$. The oriented-area identity
\[
 \operatorname{Im}\bigl(\overline{\mu_1(a)}\mu_2(a)\bigr)
       =A\det_{\R}T_a
\]
then proves the displayed formula for $\operatorname{Im}r_a$.
In particular,
\[
 \Lambda_h=c\bar\Lambda,\qquad r_h=\bar\tau,\qquad
 j(\Lambda_h)=\overline{j(E)}.
\]
The sign of $\operatorname{Im}r_h$ is negative; reversing the
orientation of the chosen lattice basis gives the usual upper
half-plane convention for its $j$-invariant.

The quasi-period law
$\zeta(u+\lambda)-\zeta(u)=\eta_\lambda$ proves that $F_a$ descends to
$E(\C)\setminus\{O\}$. Since $\zeta'=-\wp$ and $\wp=x$, it also gives
$F_a^*(dz)=-(x+a)\omega$. In a local parameter at $O$, this differential
has a double pole with nonzero leading coefficient. A holomorphic
extension of $F_a$ would pull the holomorphic differential $dz$ on the
target elliptic curve back to a holomorphic differential, which is
impossible. An algebraic map from the punctured smooth curve to this
proper elliptic curve would extend at $O$ by the valuative criterion
of properness. Therefore $F_a$ is not algebraic.

Now assume $a\in\Qb$. If $E$ is non-CM, its four numbers
$\omega_1,\omega_2,\eta_1,\eta_2$ are linearly independent over $\Qb$.
Here is the precise period-theorem input. Apply
\cite[Lemma 15.6, pp.136--138 of the opened manuscript]{EllHuberWustholz2022}
to the $1$-motive $[0\to E]$. In Definition~15.1 it is reduced and
saturated: both auxiliary lattice and torus are zero, and its
endomorphism algebra is $\operatorname{End}^0(E)=\Q$.
Its Betti basis therefore has two elements and its de Rham basis has
two elements. The toric and lattice rows in the lemma are absent;
the lemma consequently asserts independence of the four entries of
the resulting $2$ by $2$ period matrix. Changing its second row from
$(-\eta_1,-\eta_2)$ to $(\eta_1,\eta_2)$ preserves that independence.
If the quotient torus had CM, its lattice ratio $r_a$ would be
imaginary quadratic. The relation
\[
 \eta_2-r_a\eta_1-a\omega_2+a r_a\omega_1=0
\]
would then be a nonzero linear relation over $\Qb$, a contradiction.
This proves that every such quotient is non-CM when $E$ is non-CM.

Suppose instead that $E$ has CM. Then $\tau$ is imaginary quadratic
and $h$ is algebraic by A.4.2. Put
\[
 U=\omega_1,\qquad
 V=c\bar\omega_1=\frac{\pi}{U\operatorname{Im}\tau},\qquad
 b=h-a\in\Qb.
\]
The two numbers $U,V$ are linearly independent over $\Qb$.
Indeed, a dependence with both coefficients nonzero would imply
$\pi=\alpha U^2$ for $\alpha\in\Qb$, contrary to the algebraic
independence of $U$ and $\pi$ in A.4.3, proved using
\cite[Corollary 2, p.343]{EllChudnovsky1978}.
A dependence with one coefficient zero is excluded by $U,V\ne0$.
The quasi-period identity now becomes
\[
 \mu_1(a)=bU+V,\qquad
 \mu_2(a)=b\tau U+\bar\tau V.
\]
If the quotient is CM, $r_a$ is algebraic, and hence
\[
 b(\tau-r_a)U+(\bar\tau-r_a)V=0.
\]
Independence gives $r_a=\bar\tau$ and
$b(\tau-\bar\tau)=0$. Since $\tau\ne\bar\tau$, this forces
$a=h$. Conversely, $a=h$ gives $\Lambda_h=c\bar\Lambda$;
conjugation preserves the nontrivial imaginary-quadratic
endomorphism algebra, so this quotient has CM.
\end{proof}

\paragraph{The remaining quotient arithmetic.}
The proved classification concerns the CM property of the quotient.
For a non-CM $E$, its actual quotient at $a=h$ has the algebraic
$j$-invariant $\overline{j(E)}$ and is itself non-CM. The displayed
maps and period-theorem calculation do not exclude the algebraicity
of this particular $a$. The H1 gate therefore remains unchanged.

\subsection{The differential scale of an algebraic quotient}

\begin{result}{A.4.35}{Raw second-kind lattice invariants}{Proved}
Let $E/\Qb$ and $\Lambda_a$ have the conventions of A.4.34.
For every $a\in\Qb$ for which $\Lambda_a$ is a full lattice, its two
raw Weierstrass invariants $g_2(\Lambda_a),g_3(\Lambda_a)$ cannot
both be algebraic. Consequently, every $d\in\C^\times$ for which
$d^{-1}\Lambda_a$ has both Weierstrass invariants algebraic is
transcendental.
\end{result}

\begin{proof}
Suppose both raw invariants were algebraic. Then
\[
 F:\quad Y^2=4X^3-g_2(\Lambda_a)X-g_3(\Lambda_a)
\]
is an elliptic curve over $\Qb$. Its holomorphic differential
$\omega_F=dX/Y$ has full integral period module $\Lambda_a$.
On $E$, the algebraic de Rham class
\[
 \alpha=-\xi-a\omega
\]
has precisely the same integral period module, because
$\int_{\gamma_i}\xi=-\eta_i$.
Apply \cite[Corollary 2, p.14 of the opened version]{EllMovasati2008}
to $(E,\alpha)$ and $(F,\omega_F)$.
Both varieties are simple, both classes are nonzero, and their full
$\Z$-period modules agree. The corollary therefore gives an algebraic
isomorphism $b:E\longrightarrow F$ satisfying
$b^*[\omega_F]=\alpha$.
The left-hand side belongs to $F^1H^1_{\mathrm{dR}}(E/\Qb)$,
since pullback preserves regular differentials. The right-hand
side has coefficient $-1$ on $\xi$ in the basis $(\omega,\xi)$,
whereas $F^1=\Qb\omega$. This is a contradiction.

For the second assertion, the exact scaling formulas are
\[
 g_2(d^{-1}\Lambda_a)=d^4g_2(\Lambda_a),\qquad
 g_3(d^{-1}\Lambda_a)=d^6g_3(\Lambda_a).
\]
If $d$ were nonzero and algebraic, algebraicity of the two scaled
invariants would imply algebraicity of both raw invariants, contrary
to the first assertion. These formulas also cover $g_2=0$ or $g_3=0$;
no division by either invariant is made.
\end{proof}

The scale in A.4.35 has a concrete interpretation for any algebraic
model of the quotient. Choose a nonzero algebraic regular differential
on that model and use its integral to identify the universal cover
with $\C$. Every origin-preserving analytic isomorphism from
$\C/\Lambda_a$ to this model then has the form
$[z]\mapsto[z/d]$, with target period lattice $d^{-1}\Lambda_a$.
Indeed its lift between universal covers is complex-linear, and its
nonzero derivative is $1/d$. The short Weierstrass equation associated
to the chosen algebraic differential has algebraic coefficients.
Thus the assertion applies to this actual comparison scale, not
merely to the existence of an abstract algebraic $j$-invariant.

\paragraph{The actual shifted point and the remaining gate.}
The same source's Theorem~2, p.3, can be tested at the precise Hodge
point. Its family and distinguished differential are
\[
 E^t:\quad Y^2=4(X-t_1)^3-t_2(X-t_1)-t_3,
 \qquad \alpha_t=X\,dX/Y.
\]
For $t=(a,g_2,g_3)$, the algebraic map
$f_a:E\to E^t$, $(x,y)\mapsto(x+a,y)$, satisfies
$f_a^*\alpha_t=\xi+a\omega$.
At $a=h$ its period row is
\[
 (-\eta_1+h\omega_1,-\eta_2+h\omega_2)
       =-c(\bar\omega_1,\bar\omega_2).
\]
Consequently every nonzero integral cycle has a nonzero period:
its value is $-c\,\overline{m\omega_1+n\omega_2}$ with
$(m,n)\ne(0,0)$, and the lattice basis is real-linearly independent.
The zero-period hypothesis of Theorem~2(2) therefore fails at this
point. Theorem~2(1) allows at most one algebraic point on a nonsingular
leaf; algebraicity of $h$ supplies this one candidate and does not
supply a second algebraic point on the same leaf.
At the actual Hodge value,
$\Lambda_h=c\bar\Lambda$ and $j(\Lambda_h)=\overline{j(E)}$.
If $h$ is algebraic, A.4.35 forces the exhibited scale $d=c$ to be
transcendental, in agreement with A.4.33. Both CM examples and the
unresolved non-CM case meet this conclusion. The statement does not
exclude algebraicity of the scale-invariant $j(\Lambda_h)$; the H1
gate remains the nonzero algebraic-$h$ question.
